% ATTENZIONE!!! 
% Per far funzionare i collegamenti ipertestuali si raccomanda di usare
%	\classname{nomedellaclasse}
% per le classi dello stesso package mentre invece 
%	\hyperref[nomedellaclasse]{\ttfamily{}nomequalificatodellaclasse}
% per le classi che non sono dello stesso package e che hanno il nome completo

% **************************************************
% Macro specifiche per il documento corrente
% **************************************************
% Nome
\newcommand{\docName}{Definizione di prodotto}
% Nome file
\newcommand{\docFileName}{definizione\_di\_prodotto.1.0.pdf}
% Versione
\newcommand{\docVers}{1.0}
% Data creazione
\newcommand{\creationDate}{2013-02-03}
% Data ultima modifica
\newcommand{\modificationDate}{2013-02-03}
% Stato in {Approvato, Non approvato}
\newcommand{\docState}{Non approvato}
% Uso in {Interno, Esterno}
\newcommand{\docUsage}{Esterno}
% Destinatari da specificare come nome1\\ &nome2\\ ecc.
\newcommand{\docDistributionList}{Prof. Tullio Vardanega\\&Prof. Riccardo Cardin\\&Dott. Gregorio Piccoli\\&Team SoftwareSynthesis}
% Redattori da specificare come nome1\\ &nome2\\ ecc.
\newcommand{\docAuthors}{Elena Zecchinato\\&Andrea Meneghinello\\&Riccardo Tresoldi}
% Approvato da
\newcommand{\approvedBy}{Andrea Rizzi}
% Verificatori
\newcommand{\verifiedBy}{Stefano Farronato\\&Marco Schivo}
% Perscorso (relativo o assoluto) che punta alla directory contenente shared/
% come sua sottodirectory (per comodità chiamiamola 'doc root').
\newcommand{\docRoot}{..}
% definire se si vuole l'indice delle tabelle
\def\INDICETABELLE{false}
% definire se si vuole l'indice delle figure
\def\INDICEFIGURE{false}

% importa il preambolo condiviso da tutti i documenti
% shared/preamble.tex
%
% Questo documento contiene la parte del preambolo condivisa e viene pertanto
% richiamato nel 'master' di tutti i documenti di progetto.  Al suo interno
% contiene le inclusioni (e le configurazioni) di tutti i package richiesti per
% la compilazione dei documenti, le macro di carattere generale e la definizione
% degli stili di pagina.

\documentclass[a4paper,10pt,openright]{article}

% **************************************************
% Macro generiche
% **************************************************
\newcommand{\team}{Software Synthesis}                    % chi siamo
\newcommand{\email}{software.synthesis@gmail.com}         % e-mail
\newcommand{\caName}{}                                    % titolo capitolato
\newcommand{\caDescr}{}                                   % descrizione
\newcommand{\inglese}[1]{\textit{#1}}

% **************************************************
% Codifica e lingua dei documenti
% **************************************************
\usepackage[utf8x]{inputenc}                              % codifica caratteri dei documenti sorgenti
\usepackage[english,italian]{babel}                       % localizzazione ai fini di sillabazione e cross-references
\usepackage[T1]{fontenc}                                  % codifica font di output

% **************************************************
% Definizione geometria della pagina
% **************************************************
\usepackage[a4paper,head=4cm,top=4.5cm,bottom=3cm,left=3cm,right=3cm,bindingoffset=5mm]{geometry}

% *************************************************
% Intestazioni e piè di pagina personalizzati
% *************************************************
\usepackage{fancyhdr}
% stile normale
\fancypagestyle{normal}{
\fancyhead{}                                              % intestazione
\fancyhead[RE,RO]{
\begin{picture}(0,0)
  \put(-410,0){\includegraphics[width=1.02\textwidth]{header_logo}}
  \put(-410,10){\sffamily\large\leftmark}
\end{picture}
\vspace{-4pt}
}
\renewcommand{\headrulewidth}{.4pt}                       % riga sotto l'intestazione
\cfoot{}                                                  % piè di pagina
\fancyfoot[RO,LE]{\sffamily
  pag.~\thepage{} di \pageref{LastPage}}                  % a dx nelle pag. dispari e a sx in quelle pari
\fancyfoot[RE,LO]{\sffamily\docFileName{} -- v.\docVers}
\renewcommand{\footrulewidth}{.4pt}                       % riga sopra il piè di pagina
}
% stile per gli indici
\fancypagestyle{toc}{
\fancyhead{}                                              % intestazione
\fancyhead[RE,RO]{
\begin{picture}(0,0)
  \put(-410,0){\includegraphics[width=1.02\textwidth]{header_logo}}
\end{picture}
}
\renewcommand{\headrule}{}                                % nessuna riga sotto l'intestazione
\cfoot{}                                                  % piè di pagina
\fancyfoot[RO,LE]{\sffamily\thepage{}}                    % a dx nelle pag. dispari e a sx in quelle pari
\fancyfoot[RE,LO]{\sffamily\docFileName{} -- v.\docVers}
\renewcommand{\footrulewidth}{.4pt}                       % riga sopra il piè di pagina
}

\pagestyle{fancy}                                         % premetto: non so usare bene le marche:
\renewcommand{\sectionmark}[1]{\markboth{#1}{#1}}         % se qualcuno ha idee migliori si faccia avanti!

% **************************************************
% Tabelle
% **************************************************
\usepackage{tabularx}                                     % tabelle di larghezza fissa con una o più colonne variabili
\usepackage{multirow}                                     % colonne con colonne che si estendono per più righe
\usepackage{booktabs}                                     % per inserire l'ambiente table e le righe orizz. nelle tabelle
\usepackage{longtable}			                          % tabelle oltre i limiti di pagina

% **************************************************
% Cross-references e collegamenti ipertestuali
% **************************************************
\usepackage[hidelinks]{hyperref}
\hypersetup{%
  colorlinks=false, linktocpage=false, pdfborder={0,0,0}, pdfstartpage=3, pdfstartview=FitV,%
  urlcolor=Cyan, linkcolor=Cyan, citecolor=Black, %pagecolor=Black,%
  pdftitle={\docName}, pdfauthor={\team}, pdfsubject={}, pdfkeywords={},%
  pdfcreator={pdflatex}, pdfproducer={pdflatex with hyperref package}%
}

% **************************************************
% Immagini e grafica
% **************************************************
\usepackage{graphicx}                                     % supporto ad aspetti avanzati delle immagini
\graphicspath{{\docRoot/pics/}}                           % percorso contenente tutti i file immagini
\usepackage{color}                                        % permette di colorare facilmente il testo

% **************************************************
% Altri pacchetti opzionali
% **************************************************     
\usepackage{lastpage}                                     % per sapere il numero totale di pagine
\usepackage{lipsum}                                       % genera "dummy text" per prove di impaginazione
\usepackage{eurosym}                                      % per il simbolo dell'euro usare \EUR{x} dove x è l'importo


% macro specifiche per il documento corrente
\newcommand{\classsection}[1]{\subsubsection{#1}\label{#1}}
\newcommand{\classname}[1]{\hyperref[#1]{\ttfamily#1}}

% Fine del preambolo e inizio del documento
\begin{document}

% Inclusione della prima pagina
% shared/firstpage.tex
%
% Questo documento definisce il contenuto della prima pagina, che si suppone
% essere uguale in tutti i documenti.  Oltre al logo e al titolo, la prima
% pagina contiene i metadati relativi al documento in cui viene inclusa.


% rimuove intestazioni e piè di pagina
\pagestyle{empty}

\begin{center}

% logo del gruppo
\includegraphics[width=1.5\textwidth]{logo}

\vspace{1in}

% titolo del documento
{\Huge\bfseries \docName}

\vspace{1in}

% tabella riepilogativa
\begin{tabularx}{.7\textwidth}{>{\bfseries\sffamily}l>{\sffamily}l}
\toprule
\multicolumn{2}{>{\sffamily}c}{Informazioni sul documento}\\
\midrule
Nome file:            & \docFileName\\
Versione:             & \docVers\\
Data creazione:       & \creationDate\\
Data ultima modifica: & \modificationDate\\
Stato:                & \docState\\
Uso:                  & \docUsage\\
Redattori:            & \docAuthors\\
Approvato da:         & \approvedBy\\
Verificatori:         & \verifiedBy\\
\bottomrule
\end{tabularx}

\end{center}

\newpage


%---------------------------RUOLI----------------------------
%FASE 1:
%Programmatori: DIEGO, STEFANO, SCHIVO
%FASE 2:
%Programmatori: MENE, TRES, ELENA

%Verificatore: SCHIVO, STEFANO, RIZZI, DIEGO (dobbiamo fare un sacco di test)
%Responsabile finale supremo: RIZZI
%------------------------------------------------------------

% Storico delle modifiche
\section*{Storia delle modifiche}
\begin{center}
\begin{longtable}{lp{.32\textwidth}lll}
\toprule
Versione & Descrizione intervento & Membro & Ruolo & Data\\
\midrule % inserire qui il contenuto della tabella
0.1 & Creazione del documento e stesura delle sezioni ``Introduzione'' e ``Riferimenti''. & &  & 2013-02-03\\
\bottomrule
\end{longtable}
\end{center}
\newpage

% inclusione dell'indice
% shared/toc.tex
%
% Questo file contiene le istruzioni che generano l'indice o gli indici del
% documento (utile nel caso in cui decidessimo di avere anche un indice delle
% tabelle e/o un indice delle figure).

\pagestyle{toc}
\pagenumbering{roman}

\tableofcontents

\newpage


% Alcuni aggiustamenti per le pagine
\pagenumbering{arabic}
\setcounter{page}{1}
\pagestyle{normal}

% Qui ha inizio il documento vero e proprio...
\newpage

\section{Introduzione}
\subsection{Scopo del prodotto}
\purpose

\subsection{Scopo del documento}
Il presente documento presenta una descrizione dettagliata dell'architettura del sistema software destinata alla realizzazione del prodotto \caName{}, coerentemente con la progettazione ad alto livello descritta nell'allegato \textit{specifica\_tecnica.2.0.pdf}.

A tal fine si riporta per ognuno dei componenti definiti nel documento di specifica tecnica una descrizione delle classi in termini di operazioni disponibili, proprietà, responsabilità e collaborazioni. Il contenuto del presente documento ha inoltre valore vincolante per i programmatori, pertanto avranno l'obbligo di attenersi alle disposizioni in esso contenute senza alcuna possibilità di deroga.

\subsection{Glossario}
\glossaryIntro

\subsection{Convenzioni di scrittura}
% vedere issue #56 al riguardo
Al fine di rendere quanto più agevole possibile la consultazione del documento da parte dei programmatori e del committente, è stata adottata una serie di accorgimenti sia a livello di riferimenti sulla nomenclatura delle classi sia a livello cromatico per campi dati e metodi. 
Tali norme possono essere consultate in dettaglio nel documento \textit{norme\_di\_progetto.3.0.pdf} allegato.
\clearpage

\section{Riferimenti}
\subsection{Normativi}
\begin{itemize}
\item[] \textit{piano\_di\_qualifica.3.0.pdf} allegato.
\item[] \textit{norme\_di\_progetto.3.0.pdf} allegato.
\item[] \textit{specifica\_tecnica.2.0.pdf} allegato
\end{itemize}

\subsection{Informativi}
\begin{itemize}
\item[] Capitolato d'appalto: \caName{}, v1.0, redatto e rilasciato dal proponente Zucchetti s.r.l. reperibile all'indirizzo \url{http://www.math.unipd.it/~tullio/IS-1/2012/Progetto/C1.pdf};
\item[] testo di consultazione: \textit{Software Engineering (8th edition) Ian Sommerville, Pearson Education | Addison Wesley};
\item[] manuale all'utilizzo dei design pattens: \textit{Design Patterns, Elementi per il riuso di software a oggetti -- (1/Ed. italiana) Eric Gamma, Richard Helm, Ralph Johnson, John Vlissides, Pearson Education};
\item[] \textit{glossario.3.0.pdf} allegato.
\end{itemize}
\clearpage

\section{Standard di progetto}

\subsection{Standard di progettazione architetturale}
Lo sviluppo del progetto ha seguito le regole architetturali specificate nel documento  \textit{norme\_di\_progetto.3.0.pdf} allegato.

\subsection{Standard di documentazione del codice}
Le regole che definiscono la documentazione del codice relativo al funzionamento del prodotto sono specificate nel documento \textit{norme\_di\_progetto.3.0.pdf} allegato.

\subsection{Standard di denominazione di entità e relazioni}
Le convenzioni relative alla denominazione delle entità e le relative relazioni sono specificate nel documento \textit{norme\_di\_progetto.3.0.pdf} allegato.

\subsection{Standard di programmazione}
Le regole relative agli standard di programmazione sono enunciate nel documento \textit{norme\_di\_progetto.3.0.pdf} allegato.

\subsection{Strumenti di lavoro}
Gli strumenti utilizzati per la stesura e lo sviluppo sono specificati nei documenti \textit{norme\_di\_progetto.3.0.pdf} e \textit{piano\_di\_qualifica.3.0.pdf} allegati.

\clearpage

\section{Specifica sotto-architettura sever}\label{sec:serverarchitecture}

\subsection{Package org.softwaresynthesis.mytalk.server}\label{sec:server}

\classsection{HibernateUtil}

\subsubsection*{Funzione}
Interfaccia usata per dichiarare che l'oggetto che la implementa ha la possibilità di essere convertito in formato \textit{json} in modo da poter essere trasferito al client.

\subsubsection*{Relazioni d'uso}

Nessuna relazione evidenziata.

\subsubsection*{Metodi}

\begin{description}
	\item{\method{+ toJson()}}\\
	Converte l'oggetto di invocazione in una stringa \textit{json} interpretabile dai client.
\end{description}

\subsection{Package org.softwaresynthesis.mytalk.server.dao}\label{sec:dao}

\classsection{HibernateUtil}

\subsubsection*{Funzione}
Inizializza un'unica \inglese{factory} per le sessioni, utilizzate da \textit{Hibernate}, per comunicare con il database.

\subsubsection*{Relazioni d'uso}

\begin{itemize}
	\item \classname{org.hibernate.SessionFactory}: necessaria per interrogare il database.
	\item \classname{org.hibernate.cfg.Configuration}: definisce i parametri necessari per la connessione con il database. Inoltre definisce i \textit{mapping} necessari tra le classi \inglese{transfer object} e le relative tabelle nel database.
\end{itemize}

\subsubsection*{Attributi}

\begin{itemize}
	\item{\memberdata{-- \underline{instance}: HibernateUtil}}\\
	Attributo usato per implementare il \inglese{pattern} singleton. Tale istanza verrà inizializzata tramite il metodo statico \method{getIstance()}, assicurando che l'attributo non sia già stato inizializzato in precedenza.
	\item{\memberdata{-- \underline{sessionFactory}: SessionFactory}}\\
	Attributo contenente la \inglese{factory} delle sessioni verso il database.
\end{itemize}

\subsubsection*{Metodi}

\begin{description}
	\item{\method{-- HibernateUtil()}}\\
	Costruttore privato della classe, definito \inglese{private} in correlazione all'applicazione del \inglese{pattern} singleton.
	\item{\method{+ getInstance(): HibernateUtil}}\\
	Metodo pubblico che ritorna l'istanza \memberdata{istance}. Il metodo controlla se \memberdata{istance} è già stata inizializzata, nel caso in cui non lo sia il metodo dovrà generare un istanza di HibernateUtil richiamando il costruttore privato \method{HibernateUtil()} e assegnare il valore ritornato all'attributo \memberdata{istance}. Il programma termina restituendo \memberdata{istance}.
	\item{\method{+ getSessionFactory(): SessionFactory}}\\
	Metodo che ritorna l'attributo \memberdata{sessionFactory}.
\end{description}

\classsection{UserDataDAO}

\subsubsection*{Funzione}
Classe che implementa il pattern DAO. Definisce le procedure d'inserimento, eliminazione ed aggiornamento di \inglese{entry} inerenti alla tabella UserData. La classe inoltre fornisce metodi per interrogare il database ed ottenere i dati utente con i quali costruire istanze di oggetti di tipo \classname{abook.IUserData}.

\subsubsection*{Relazioni d'uso}

\begin{itemize}
	\item \texttt{java.util.List}: usata per memorizzare e ritornare i dati restituiti dalle interrogazioni a database.
	\item \classname{org.hibernate.Query}: classe rappresentante le \textit{query} d'interrogazione verso il database
	\item \texttt{org.hibernate.Session}: classe che rappresenta la sessione di ``lavoro'' verso il database.
	\item \texttt{org.hibernate.SessionFactory}: classe avente compito di instanziare nuove sessioni.
	\item \texttt{org.hibernate.Transaction}: classe che rappresenta una transazione del database.
	\item \classname{abook.IUserData}: interfaccia del package \textit{abook} usata per definire un generico utente. La classe può ritornare istanze di tipo \classname{abook.IUserData} come risposta a delle \textit{query} d'interrogazione al database.
	
\end{itemize}

\subsubsection*{Attributi}

Nessun attributo evidenziato.

\subsubsection*{Metodi}

\begin{description}

	\item{\method{+ insert(user: IUserData): boolean}}\\
	Metodo usato per inserire un utente registrato nel sistema, riceve come parametro un istanza di un oggetto \classname{abook.IUserData}. Il metodo dovrà eseguire nell'ordine le seguenti operazioni:
		\begin{itemize}
			\item[1)] ottenere dalla classe \classname{HibernateUtil} l'istanza tramite il metodo statico \method{getIstance()};
			\item[2)] ottenere dall'istanza di tipo \classname{HibernateUtil} (ottenuta al punto precedente), un oggetto di tipo \texttt{org.hibernate.SessionFactory}, tramite il metodo \method{getSessionFactory()};
			\item[3)] aprire la sessione ottenuta al punto precedente;
			\item[4)] avviare una transizione a partire dall'istanza \texttt{org.hibernate.SessionFactory} del punto 2. Tale istanza va memorizzata in un opportuno attributo (interno al metodo stesso) di tipo \texttt{org.hibernate.Transaction};
			\item[5)] eseguire il comando \inglese{save} a partire dalla sessione aperta e passando come parametro l'attributo \memberdata{user}.
			\item[6)] eseguire il comando \inglese{commit} per confermare l'operazione del punto precedente.
		\end{itemize}
		
Poiché l'iter presentato potrà causare il lancio di un'eccezione, il codice dovrà essere gestito in un blocco \inglese{try-catch}. Nel caso incorra tale evento, il flusso di processo dovrà ristabilire la situazione iniziale richiamando il metodo \texttt{rollback()} dell'istanza di tipo \texttt{org.hibernate.Transaction} creata al punto 4 dell'iter. In ogni caso il metodo deve terminare chiudendo la sessione aperta al punto 2.

	\item{\method{+ delete(user: IUserData): boolean}}\\
	Metodo usato per eliminare un utente registrato nel sistema, riceve come parametro un istanza di un oggetto \classname{abook.IUserData}. Il metodo dovrà eseguire nell'ordine le seguenti operazioni:
		\begin{itemize}
			\item[1)] ottenere dalla classe \classname{HibernateUtil} l'istanza tramite il metodo statico \method{getIstance()};
			\item[2)] ottenere dall'istanza di tipo \classname{HibernateUtil}(ottenuta al punto precedente), un oggetto di tipo \texttt{org.hibernate.SessionFactory}, tramite il metodo \method{getSessionFactory()};
			\item[3)] aprire la sessione ottenuta al punto precedente;
			\item[4)] avviare una transizione a partire dall'istanza \texttt{org.hibernate.SessionFactory} del punto 2. Tale istanza va memorizzato in un opportuno attributo (interno al metodo stesso) di tipo \texttt{org.hibernate.Transaction};
			\item[5)] eseguire il comando \inglese{delete} a partire dalla sessione aperta e passando come parametro l'attributo \memberdata{user}.
			\item[6)] eseguire il comando \inglese{commit} per confermare l'operazione del punto precedente.
		\end{itemize}
		
Poiché l'iter presentato può causare il lancio di un'eccezione, il codice dovrà essere gestito in un blocco \inglese{try-catch}. Nel caso incorra tale evento, il flusso di processo dovrà ristabilire la situazione iniziale richiamando il metodo \texttt{rollback()} dell'istanza di tipo \texttt{org.hibernate.Transaction} creata al punto 4 dell'iter. In ogni caso il metodo deve terminare chiudendo la sessione aperta al punto 2.

\item{\method{+ update(user: IUserData): boolean}}\\
	Metodo usato per modificare i dati di un utente registrato nel sistema, riceve come parametro un istanza di un oggetto \classname{abook.IUserData}. Il metodo dovrà eseguire nell'ordine le seguenti operazioni:
		\begin{itemize}
			\item[1)] ottenere dalla classe \classname{HibernateUtil} l'istanza tramite il metodo statico \method{getIstance()};
			\item[2)] ottenere dall'istanza di tipo \classname{HibernateUtil}(ottenuta al punto precedente), un oggetto di tipo \texttt{org.hibernate.SessionFactory}, tramite il metodo \method{getSessionFactory()};
			\item[3)] aprire la sessione ottenuta al punto precedente;
			\item[4)] avviare una transizione a partire dall'istanza \texttt{org.hibernate.SessionFactory} del punto 2. Tale istanza va memorizzato in un opportuno attributo (interno al metodo stesso) di tipo \texttt{org.hibernate.Transaction};
			\item[5)] eseguire il comando \inglese{update} a partire dalla sessione aperta e passando come parametro l'attributo \memberdata{user}.
			\item[6)] eseguire il comando \inglese{commit} per confermare l'operazione del punto precedente.
		\end{itemize}
		
Poiché l'iter presentato può causare il lancio di un'eccezione, il codice dovrà essere gestito in un blocco \inglese{try-catch}. Nel caso incorra il lancio tale evento, il flusso di processo dovrà ristabilire la situazione iniziale richiamando il metodo \texttt{rollback()} dell'istanza di tipo \texttt{org.hibernate.Transaction} creata al punto 4 dell'iter. In ogni caso il metodo deve terminare chiudendo la sessione aperta al punto 2.

	\item{\method{+ getByNameAndSurname(name: String, surname: String): List<IUserData>}}\\
	Metodo usato per effettuare un interrogazione al database atta ad ottenere la lista degli utenti aventi un determinato nome e cognome, riceve come parametri il nome e il cognome da ricercare. Il metodo dovrà eseguire nell'ordine le seguenti operazioni:
		\begin{itemize}
			\item[1)] ottenere dalla classe \classname{HibernateUtil} l'istanza tramite il metodo statico \method{getIstance()};
			\item[2)] ottenere dall'istanza di tipo \classname{HibernateUtil}(ottenuta al punto precedente), un oggetto di tipo \texttt{org.hibernate.SessionFactory}, tramite il metodo \method{getSessionFactory()};
			\item[3)] aprire la sessione ottenuta al punto precedente;
			\item[4)] eseguire il comando \inglese{createQuery()} a partire dall'istanza che definisce la sessione, passando come parametro una stringa rappresentante la \inglese{query} d'interrogazione al database (la \inglese{query} deve restituire dati validi ad essere passati a tale metodo). Una prima analisi evidenzia che la \inglese{query} corretta è: 
\begin{verbatim}
from UserData as u where u.name = :name or u.surname = :surname
\end{verbatim}
			\item[5)] eseguire una serie di comandi di tipo \texttt{setString} per ogni elemento della \inglese{query} d'interrogazione (dove per elemento si intende ogni valore di filtraggio presente nelle clausole \inglese{where} delle \inglese{query} HQL, segnato come \texttt{:nome\_variabile}).
			\item[6)] eseguire il comando \inglese{commit} per l'interrogare il database ed ottenere la lista di valori da ritornare.
		\end{itemize}
		
Il metodo deve terminare chiudendo la sessione aperta al punto 2.
	
	\item{\method{+ getByEmail(mail: String): IUserData}}\\
	Metodo usato per eseguire un interrogazione al database atta ad ottenere l'utente avente un determinato indirizzo e-mail, esso riceve come parametro l'indirizzo e-mail da ricercare. Il metodo dovrà eseguire nell'ordine le seguenti operazioni:
		\begin{itemize}
			\item[1)] ottenere dalla classe \classname{HibernateUtil} l'istanza tramite il metodo statico \method{getIstance()};
			\item[2)] ottenere dall'istanza di tipo \classname{HibernateUtil}(ottenuta al punto precedente), un oggetto di tipo \texttt{org.hibernate.SessionFactory}, tramite il metodo \method{getSessionFactory()};
			\item[3)] aprire la sessione ottenuta al punto precedente;
			\item[4)] eseguire il comando \inglese{createQuery()} a partire dall'istanza che definisce la sessione, passando come parametro una stringa rappresentante la \inglese{query} d'interrogazione al database (la \inglese{query} deve restituire dati validi ad essere passati a tale metodo). Una prima analisi evidenzia che la \inglese{query} corretta è: 
\begin{verbatim}
from UserData as u where u.mail = :mail
\end{verbatim}
			\item[5)] eseguire una serie di comandi di tipo \texttt{setString} per ogni elemento della \inglese{query} d'interrogazione (per elemento si intende ogni valore di filtraggio presente nelle clausole \inglese{where} delle \inglese{query} HQL, segnato come \texttt{:nome\_variabile}).
			\item[6)] eseguire il comando \inglese{commit} per eseguire l'interrogazione a database ed ottenere la lista di valori da ritornare.
		\end{itemize}
		
Il metodo deve terminare chiudendo la sessione aperta al punto 2.	
	
	
	\item{\method{+ searchGeneric(value: String): List<IUserData>}}\\
		Metodo usato per eseguire un interrogazione generica al database atta ad ottenere una lista di utenti aventi: o un determinato indirizzo e-mail, o un determinato nome o un determinato cognome; riceve come parametro la stringa da ricercare e dovrà eseguire nell'ordine le seguenti operazioni:
		\begin{itemize}
			\item[1)] ottenere dalla classe \classname{HibernateUtil} l'istanza tramite il metodo statico \method{getIstance()};
			\item[2)] ottenere dall'istanza di tipo \classname{HibernateUtil}(ottenuta al punto precedente), un oggetto di tipo \texttt{org.hibernate.SessionFactory}, tramite il metodo \method{getSessionFactory()};
			\item[3)] aprire la sessione ottenuta al punto precedente;
			\item[4)] eseguire il comando \inglese{createQuery()} a partire dall'istanza che definisce la sessione, passando come parametro una stringa rappresentante la \inglese{query} d'interrogazione al database (la \inglese{query} deve restituire dati validi ad essere passati a tale metodo). Una prima analisi evidenzia che la \inglese{query} corretta è: 
\begin{verbatim}
from UserData as u where u.mail like :mail or u.name like :name or
\end{verbatim}
\begin{verbatim}
u.surname like :surname
\end{verbatim}
			\item[5)] eseguire una serie di comandi di tipo \texttt{setString} per ogni elemento della \inglese{query} d'interrogazione (per elemento si intende ogni valore di filtraggio presente nelle clausole \inglese{where} delle \inglese{query} HQL, segnato come \texttt{:nome\_variabile}).
			\item[6)] eseguire il comando \inglese{commit} per eseguire l'interrogazione a database ed ottenere la lista di valori da ritornare.
		\end{itemize}
		
Il metodo deve terminare chiudendo la sessione aperta al punto 2.

\end{description}

\subsection{Package org.softwaresynthesis.mytalk.server.abook}\label{sec:abook}

\classsection{IUserData}

\subsubsection*{Funzione}
Interfaccia rappresentante il comportamento di un generico utente del sistema. L'interfaccia dovrà definire dei metodi di tipo get e set per i dati d'interesse.

\subsubsection*{Relazioni d'uso}

\begin{itemize}
	\item \texttt{org.softwaresynthesis.mytalk.server.IMyTalkObject}: interfaccia da estendere. Ogni oggetto che implementerà l'interfaccia \classname{IUserData} dovrà essere in grado di convertire il proprio contenuto informativo in formato Json.
	\item \classname{AddressBookEntry}: l'interfaccia definisce dei metodi per la manipolazione di dati \classname{AddressBookEntry}.
\end{itemize}

\subsubsection*{Metodi}
\begin{description}
	\item{\method{+ getId(): Long}}\\
	Restituisce l'identificatore univoco di uno \classname{IUserData}.
	\item{\method{+ getEmail(): String}}\\
	Restituisce l'indirizzo e-mail con cui uno \classname{IUserData} si è registrato nel sistema \caName.
	\item{\method{+ setEmail(mail: String): void}}\\ 
	Imposta l'indirizzo e-mail con cui si registra nel sistema \caName uno \classname{IUserData}.
	\item{\method{+ getPassword(): String}}\\
	Restituisce la password di accesso al sistema \caName di uno \classname{IUserData}.
	\item{\method{+ setPassword(password: String): void}}\\
	Imposta la password di accesso al sistema di uno \classname{IUserData}.
	\item{\method{+ getQuestion(): String}}\\
	Restituisce la domanda segreta, scelta da uno \classname{IUserData}, per il recupero della \textit{password} smarrita di accesso al sistema \caName.
	\item{\method{+ setQuestion(question: String): void}}\\
	Imposta la domanda segreta, scelta da uno \classname{IUserData}, per il recupero della \textit{password} smarrita di accesso al sistema \caName.
	\item{\method{+ getAnswer(): String}}\\
	Restituisce la risposta alla domanda per il recupero della \textit{password} smarrita di accesso al sistema \caName.
	\item{\method{+ setAnswer(answer: String): void}}\\
	Imposta la risposta alla domanda segreta per il recupero della \textit{password} di accesso al sistema \caName.
	\item{\method{+ getName(): String}}\\
	Restituisce il nome di uno \classname{IUserData}.
	\item{\method{+ setName(name: String): void}}\\
	Imposta il nome di uno \classname{IUserData}.
	\item{\method{+ getSurname(): String}}\\
	Restituisce il cognome di uno \classname{IUserData}.
	\item{\method{+ setSurname(surname: String): void}}\\
	Imposta il cognome di uno \classname{IUserData}.
	\item{\method{+ getPicturePath(): String}}\\
	Restituisce una stringa con il percorso dell'immagine del profilo di uno \classname{IUserData}.
	\item{\method{+ setPicturePath(path: String): void}}\\
	Imposta il percorso dell'immagine profilo di uno \classname{IUserData}.
	\item{\method{+ getAddressBook(): Set<AddressBookEntry>}}\\
	Metodo che ritorna il la rubrica dell'utente sotto forma d'insieme di \classname{AddressBookEntry}.
	\item{\method{+ addAddressBookEntry(entry: AddressBookEntry): void}}\\
	Metodo usato per aggiungere una nuova \classname{AddressBookEntry} all'insieme di \inglese{entry} che costituisce la rubrica utente.
\end{description}

\classsection{IGroup}

\subsubsection*{Funzione}
Interfaccia rappresentante un gruppo di una rubrica utente del sistema \caName.

\subsubsection*{Relazioni d'uso}

\begin{itemize}
		\item \texttt{org.softwaresynthesis.mytalk.server.IMyTalkObject}: interfaccia da estendere. Ogni oggetto che implementerà l'interfaccia \classname{IGroup} dovrà essere in grado di convertire il proprio contenuto informativo in formato \textit{Json}.
\end{itemize}

\subsubsection*{Metodi}
\begin{description}
	\item{\method{+ getId(): Long}}\\
	Restituisce l'identificativo univoco di uno gruppo di una rubrica utente.
	\item{\method{+ getName(): String}}\\
	Restituisce il nome di un gruppo di una rubrica utente.
	\item{\method{+ setName(name: String): void}}\\
	Imposta il nome di un gruppo di una rubrica utente.
\end{description}

\classsection{IAddressBookEntry}

\subsubsection*{Funzione}
Interfaccia rappresentante una \textit{entry} di una rubrica utente del sistema mytalk.

\subsubsection*{Relazioni d'uso}
\begin{itemize}
	\item \texttt{org.softwaresynthesis.mytalk.server.IMyTalkObject}: interfaccia da estendere. Ogni oggetto che implementerà l'interfaccia \classname{IAddressBookEntry} dovrà essere in grado di convertire il proprio contenuto informativo in formato \textit{Json}.
	\item \classname{IUserData}: l'interfaccia \classname{IAddressBookEntry} definisce più metodi che restituiscono oggetti aventi tipo di ritorno \classname{IUserData}, essi sono i metodi \inglese{get} per ottenere il ``proprietario'' della rubrica e per ottenere l'utente registrato nella rubrica. Analogamente \classname{IUserData} viene usato come parametro d'ingresso per i metodi \inglese{set} collegati ai metodi già citati.
\end{itemize}

\subsubsection*{Metodi}
\begin{description}
	\item{\method{+ getId(): Long}}\\
	Restituisce l'identificativo univoco di una \inglese{entry} di una rubrica utente del sistema \caName.
	\item{\method{+ getEntry(): IUserData}}\\
	Restituisce un istanza di un oggetto avente tipo \classname{IUserData} rappresentante un contatto della rubrica.
	\item{\method{+ setEntry(contact: IUserData): void}}\\
	Imposta l'utente \classname{IUserData} (passato come parametro d'ingresso) come contatto della rubrica.
	\item{\method{+ getGroup(): IGroup}}\\
	Restituisce il gruppo a cui appartiene lo \classname{IUserData} registrato nella rubrica.
	\item{\method{+ setGroup(group: IGroup): void}}\\
	Imposta il gruppo di appartenenza dello \classname{IUserData} registrato nella rubrica.
	\item{\method{+ getOwner(): IUserData}}\\
	Restituisce lo \classname{IUserData} possessore dell'\inglese{entry} corrente della rubrica
	\item{\method{+ setOwner(owner: IUserData ): void}}\\
	Imposta l'utente \classname{IUserData} possessore della \inglese{entry} della rubrica.
\end{description}

\classsection{UserData}

\subsubsection*{Funzione}
Implementazione dell'interfaccia \classname{IUserData}. Un istanza della classe dovrà rappresentare un generico utente del sistema definendone gli attributi e i metodi per impostare ed ottenere il contenuto dei medesimi.

\subsubsection*{Relazioni d'uso}

\begin{itemize}
	\item{IUserData}: interfaccia da implementare.
\end{itemize}

\subsubsection*{Attributi}

\begin{itemize}
	\item{\memberdata{-- id: long}}
	Attributo che definisce il codice identificativo con il quale l'utente è registrato nel database del sistema.
	\item{\memberdata{-- mail: String}}
	Attributo che definisce l'indirizzo e-mail con il quale l'utente si è registrato nel sistema.
	\item{\memberdata{-- password: String}}
	Attributo che definisce la password per il \inglese{login} dell'utente nel sistema.
	\item{\memberdata{-- question: String}}
	Attributo che definisce la domanda segreta usata dall'utente in caso di smarrimento della password.
	\item{\memberdata{-- answer: String}}
	Attributo che definisce la risposta alla domanda segreta definita nell'attributo \memberdata{question}.
	\item{\memberdata{-- name: String}}
	Attributo che definisce il nome dell'utente.
	\item{\memberdata{-- surname: String}}
	Attributo che definisce il cognome dell'utente.
	\item{\memberdata{-- path: String}}
	Attributo che definisce il percorso (sul server) in cui è memorizzata l'immagine del profilo dell'utente.
	\item{\memberdata{-- addressBook: Set<AddressBookEntry>}}
	Attributo che definisce l'insieme di \classname{AddressBookEntry} che costituiscono la rubrica dell'utente.
	
\end{itemize}

\subsubsection*{Metodi}

\begin{description}
	\item{\method{+ getId(): Long}}\\
	Restituisce l'identificatore univoco di un utente, ritornando l'attributo \memberdata{id}.
	\item{\method{\# setId(id: long): void}}\\
	Imposta l'indirizzo \textit{id} con cui l'utente si registra nel sistema \caName. Il metodo sovrascrive il contenuto dell'attributo \memberdata{id} con il valore tipo \texttt{long} ricevuto come parametro d'ingresso. ritornando l'attributo \memberdata{id}.
	\item{\method{+ getEmail(): String}}\\
	Restituisce l'indirizzo e-mail con cui uno l'utente si è registrato nel sistema \caName, ritornando il contenuto dell'attributo \memberdata{mail}.
	\item{\method{+ setEmail(mail: String): void}}\\ 
	Imposta l'indirizzo e-mail con cui l'utente si registra nel sistema \caName. Il metodo non fa altro che sovrascrivere il contenuto dell'attributo \memberdata{mail} con il valore tipo \texttt{String} ricevuto come parametro d'ingresso.
	\item{\method{+ getPassword(): String}}\\
	Restituisce la password dell'utente, ritornando il valore contenuto nell'attributo \memberdata{password}.
	\item{\method{+ setPassword(password: String): void}}\\
	Imposta la password di accesso al sistema, sovrascrivendo il contenuto dell'attributo \memberdata{password} con il valore di tipo \texttt{String} ricevuto come parametro d'ingresso.
	\item{\method{+ getQuestion(): String}}\\
	Restituisce la domanda segreta, scelta dall'utente, per il recupero della password di accesso al sistema \caName. Nello specifico il metodo restituisce il contenuto dell'attributo \memberdata{question}.
	\item{\method{+ setQuestion(question: String): void}}\\
	Imposta la domanda segreta da inserire in caso di smarrimento della password. Il metodo sovrascrive il contenuto dell'attributo \memberdata{question} con il valore tipo \texttt{String} ricevuto come parametro d'ingresso
	\item{\method{+ getAnswer(): String}}\\
	Restituisce la risposta alla domanda per il recupero della password (smarrita) di accesso al sistema \caName. Il metodo ritorna il contenuto dell'attributo \memberdata{answer}.
	\item{\method{+ setAnswer(answer: String): void}}\\
	Imposta la risposta alla domanda segreta per il recupero della password. Il metodo sovrascrive il contenuto dell'attributo \memberdata{answer} con il valore tipo \texttt{String} passato come parametro d'ingresso.
	\item{\method{+ getName(): String}}\\
	Restituisce il nome dell'utente ritornando il contenuto dell'attributo \memberdata{name}.
	\item{\method{+ setName(name: String): void}}\\
	Imposta il nome dell'utente sovrascrivendo il contenuto dell'attributo \memberdata{name} con il valore tipo \texttt{String} passato al metodo come parametro d'ingresso.
	\item{\method{+ getSurname(): String}}\\
	Restituisce il cognome dell'utente restituendo il contenuto dell'attributo \memberdata{surname}.
	\item{\method{+ setSurname(surname: String): void}}\\
	Imposta il cognome dell'utente sovrascrivendo il contenuto dell'attributo \memberdata{surnamename} con il valore tipo \texttt{String} passato al metodo come parametro d'ingresso.
	\item{\method{+ getPicturePath(): String}}\\
	Restituisce una stringa con il percorso dell'immagine del profilo dell'utente, restituendo il contenuto dell'attributo \memberdata{path}.
	\item{\method{+ setPicturePath(path: String): void}}\\
	Imposta il percorso dell'immagine profilo di un utente, sovrascrivendo il contenuto dell'attributo \memberdata{path} con il valore tipo \texttt{String} passato al metodo come parametro d'ingresso.
	\item{\method{+ getState(): State}}\\
	Restituisce lo stato in cui si trova l'utente, ritornando il contenuto dell'attributo \memberdata{state}.
	\item{\method{+ setState(state: State): void}}\\
	Imposta lo stato in cui si trova l'utente, sovrascrivendo il contenuto dell'attributo \memberdata{state} con il valore ricevuto come parametro d'ingresso.
	\item{\method{+ getAddressBook(): Set<AddressBookEntry>}}\\
	Metodo che ritorna il contenuto di \memberdata{addressBook}.
	\item{\method{+ addAddressBookEntry(entry: AddressBookEntry): void}}\\
	Metodo usato per aggiungere ad \memberdata{addressBook} una nuova \classname{AddressBookEntry} passata come parametro d'ingresso.
	\item{\method{+ toJson(): String}}\\
	Metodo usato per ritornare il contenuto di un istanza di \classname{UserData} sotto forma di stringa formattata in \textit{Json}. La stringa ritornata deve corrispondere al seguente formato:\\\\
	\verb|{name:"mio_nome",surname:"mio_cognome",email:"mia_mail"|\\\verb|,picturePath:"mia_immagine",id:"mio_id"}|\\
	
	dove i valori tra virgolette rappresentano il contenuto dei rispettivi campi dati contenuti nella classe.
\end{description}

\classsection{Group}

\subsubsection*{Funzione}
Implementazione dell'interfaccia \classname{IGroup}.

\subsubsection*{Relazioni d'uso}

\begin{itemize}
	\item \classname{IGroup}: interfaccia d'implementazione.
\end{itemize}

\subsubsection*{Attributi}

\begin{itemize}
	\item{\memberdata{-- id: long}}
	Attributo del codice identificativo del gruppo.
	\item{\memberdata{-- name: String}}:
	Attributo del nome del gruppo.
\end{itemize}

\subsubsection*{Metodi}

\begin{description}
	\item{\method{+ getId(): Long}}\\
	Restituisce l'identificativo univoco di uno gruppo di una rubrica utente.
	\item{\method{+ getName(): String}}\\
	Restituisce il nome di un gruppo di una rubrica utente.
	\item{\method{+ setName(name: String): void}}\\
	Imposta il nome di un gruppo di una rubrica utente.
	\item{\method{+ toJson(): String}}\\
	Metodo usato per ritornare il contenuto di un istanza di \classname{Group} sotto forma di stringa formattata in \textit{Json}. La stringa ritornata deve corrispondere al seguente formato:\\\\
	\verb|{id:"mio_id",name:"mio_nome"}|\\
	
	dove i valori tra virgolette rappresentano il contenuto dei rispettivi campi dati contenuti nella classe.
\end{description}

\classsection{AddressBookEntry}

\subsubsection*{Funzione}
Implementazione dell'interfaccia IAddressBookEntry.

\subsubsection*{Relazioni d'uso}

\begin{itemize}
	\item \classname{IAddressBookEntry}: interfaccia d'implementazione della classe.
	\item \classname{IUserData}: usata per definire gli attributi destinati a identificare il possessore dell'istanza di \classname{AddressBookEntry} e il relativo contatto in essa registrato.
\end{itemize}

\subsubsection*{Attributi}

\begin{itemize}
	\item{\memberdata{-- id: long}}
	Attributo del codice identificativo della classe.
	\item{\memberdata{-- group: IGroup}}
	Attributo destinato a identificare il gruppo a cui appartiene il contatto \classname{IUserName} registrato nella classe. Si ricorda che il contatto può anche non appartenere ad alcun gruppo.
	\item{\memberdata{-- contact: IUserData}}
	Attributo destinato ad identificare il contatto registrato nell'istanza di \classname{AddressBookEntry}.
	\item{\memberdata{-- owner: IUserData}}:
	Attributo destinato ad identificare il possessore dell'istanza di \classname{AddressBookEntry}.
	\item{\memberdata{-- blocked: boolean}}
	Attributo booleano necessario per bloccare il contatto. Tale blocco avviene impostando l'attributo a \texttt{true}.
\end{itemize}


\subsubsection*{Metodi}

\begin{description}
	\item{\method{+ getId(): Long}}\\
	Restituisce l'identificativo univoco della \inglese{entry} di una rubrica utente del sistema \caName, nello specifico il contenuto dell'attributo \memberdata{id}.
	\item{\method{+ getEntry(): IUserData}}\\
	Restituisce il contenuto dell'attributo \memberdata{contact}.
	\item{\method{+ setEntry(contact: IUserData): void}}\\
	Imposta l'utente \classname{IUserData} (passato come parametro d'ingresso) come contatto della rubrica, nello specifico il contenuto dell'attributo \memberdata{contact}.
	\item{\method{+ getGroup(): IGroup}}\\
	Restituisce il gruppo a cui appartiene lo \classname{IUserData} registrato nella rubrica, nello specifico il contenuto dell'attributo \memberdata{group}.
	\item{\method{+ setGroup(group: IGroup): void}}\\
	Imposta il gruppo di appartenenza dello \classname{IUserData} registrato nella rubrica, nello specifico il contenuto dell'attributo \memberdata{group}.
	\item{\method{+ getOwner(): IUserData}}\\
	Restituisce lo \classname{IUserData} possesore di questa \inglese{entry} della rubrica, nello specifico il contenuto dell'attributo \memberdata{owner}.
	\item{\method{+ setOwner(owner: IUserData ): void}}\\
	Imposta l'utente \classname{IUserData} possessore della \textit{entry} della rubrica, nello specifico il contenuto dell'attributo \memberdata{owner}.
	\item{\method{+ toJson(): String}}\\
	Metodo usato per ritornare il contenuto di un istanza di \classname{AddressBookEntry} sotto forma di stringa formattata in \textit{Json}. La stringa ritornata deve corrispondere al seguente formato:\\\\
	\verb|{id:"mio_id",contact:"mio_contatto",group:"mia_gruppo",blocked:"bloccato"}|\\
	
	dove i valori tra virgolette rappresentano il contenuto dei rispettivi campi dati contenuti nella classe.

\end{description}

\subsection{Package org.softwaresynthesis.mytalk.server.call}\label{sec:call}

\classsection{ICall}

\subsubsection*{Funzione}
Interfaccia che rappresenta una chiamata effettuata dal sistema \caName. Gli oggetti che implementano tale interfaccia vengono usati per rappresentare lo storico delle chiamate di un utente.

\subsubsection*{Relazioni d'uso}
\begin{itemize}
	\item \texttt{java.util.Date}: tipo utilizzato per definire la data d'inizio e fine di una chiamata.
\end{itemize}

\subsubsection*{Metodi}
\begin{description}
	\item{\method{+ getId(): Long}}\\
	Restituisce l'identificativo univoco della chiamata.
	\item{\method{+ getStartDate(): Date}}\\
	Restituisce un istanza di \texttt{java.util.Date} di avvio della chiamata.
	\item{\method{+ setStartDate(startDate: Date): void}}\\
	Imposta la data di avvio della chiamata.
	\item{\method{+ getEndDate(): Date}}\\
	Restituisce la data in cui termina la chiamata
	\item{\method{+ setEndDate(endDate: Date): void}}\\
	Imposta la data in cui termina la chiamata.
\end{description}

\classsection{Call}

\subsubsection*{Funzione}
Classe che implementa l'interfaccia \classname{ICall}.

\subsubsection*{Relazioni d'uso}
\begin{itemize}
	\item \texttt{java.util.Date}: tipo utilizzato per definire la data d'inizio e fine di una chiamata.
\end{itemize}

\subsubsection*{Attributi}

\begin{itemize}
	\item{\memberdata{-- id: long}}
	Attributo contenente il codice identificativo della chiamata.
	\item{\memberdata{-- startDate: Date}}
	Attributo di tipo \texttt{java.util.Date} contenente la data (compresa l'ora) d'inizio della chiamata.
	\item{\memberdata{-- endDate: Date}}
	Attributo di tipo \texttt{java.util.Date} contenente la data (compresa l'ora) in cui la chiamata è terminata.
\end{itemize}

\subsubsection*{Metodi}
\begin{description}
	\item{\method{+ getId(): Long}}\\
	Restituisce l'identificativo univoco della chiamata ritornando il contenuto dell'attributo \memberdata{id}.
	\item{\method{+ getStartDate(): Date}}\\
	Restituisce il contenuto del campo \memberdata{startDate}.
	\item{\method{+ setStartDate(startDate: Date): void}}\\
	Imposta la data di avvio della chiamata, sovrascrivendo il contenuto \memberdata{startData}.
	\item{\method{+ getEndDate(): Date}}\\
	Restituisce il contenuto del campo \memberdata{endDate}.
	\item{\method{+ setEndDate(endDate: Date): void}}\\
	Imposta la data di avvio della chiamata, sovrascrivendo il contenuto \memberdata{endData}.
\end{description}

\subsection{Package org.softwaresynthesis.mytalk.server.message}\label{sec:message}

\classsection{IMessage}

\subsubsection*{Funzione}
Interfaccia che rappresenta un messaggio di segreteria del sistema \caName.

\subsubsection*{Relazioni d'uso}
\begin{itemize}
	\item \texttt{java.util.Date}: tipo utilizzato per definire la data in cui è stato registrato un messaggio.
	\item \classname{abook.IUserData}: usata per rappresentare il mittente e il destinatario del messaggio di segreteria.
\end{itemize}

\subsubsection*{Metodi}
\begin{description}
	\item{\method{+ getId(): Long}}\\
	Restituisce l'identificativo univoco del messaggio.
	\item{\method{+ getSender(): IUserData}}\\
	Restituisce un istanza di tipo classname{abook.IUserData} che rappresenta il mittente del messaggio.
	\item{\method{+ setSender(sender: IUserData): void}}\\
	Imposta il mittente del messaggio.
	\item{\method{+ getReceiver(): IUserData}}\\
	Restituisce un istanza di tipo classname{abook.IUserData} che rappresenta il destinatario del messaggio.
	\item{\method{+ setReceiver(receiver: IUserData): void}}\\
	Imposta il destinatario del messaggio.
	\item{\method{+ isNew(): boolean}}\\
	Restituisce un valore booleano che identifica lo stato del messaggio (``già letto'' se il valore ritornato è \texttt{true}, ``da leggere'' se il valore ritornato è \texttt{false}).
	\item{\method{+ setNew(status: boolean): void}}\\
	Imposta il messaggio come ``già ascoltato'' o come ``da ascoltare''.
	\item{\method{+ isVideo(): boolean}}\\
	Restituisce un booleano che determina se si tratta di un messaggio audio oppure audio-video.
	\item{\method{+ setVideo(video: boolean): void }}\\	
	Imposta il messaggio come messaggio audio-video o come messaggio contenente solamente una traccia audio.
	\item{\method{+ getDate(): Date}}\\
	Restituisce la data in cui il mittente ha lasciato il messaggio nella segreteria del destinatario.
	\item{\method{+ setDate(date: Date): void}}\\
	Imposta la data in cui il mittente ha lasciato il messaggio nella segreteria del destinatario.
\end{description}

\classsection{Message}

\subsubsection*{Funzione}
Classe che implementa l'interfaccia \classname{IMessage}.

\subsubsection*{Relazioni d'uso}
\begin{itemize}
	\item \texttt{java.util.Date}: tipo utilizzato per definire la data in cui è stato registrato un messaggio.
	\item \classname{abook.IUserData}: usata per rappresentare il mittente e il destinatario del messaggio di segreteria.
\end{itemize}

\subsubsection*{Attributi}

\begin{itemize}
	\item{\memberdata{-- id: long}}
	Attributo contenente il codice identificativo del messaggio.
	\item{\memberdata{-- sender: IUserData}}
	Attributo contenente lo \classname{IUserData} che rappresenta il destinatario del messaggio.
	\item{\memberdata{-- receiver: IUserData}}
	Attributo contenente lo \classname{IUserData} che rappresenta il destinatario del messaggio.
	\item{\memberdata{-- status: boolean}}
	Attributo contenente un valore booleano che identifica se il messaggio è stato visionato/ascoltato o meno. L'attributo se impostato a \texttt{true}, designa il messaggio come \textit{ascoltato}. Invece se impostato a \texttt{false} identifica il messaggio come \textit{ancora da ascoltare}.
	\item{\memberdata{-- video: boolean}}
	Attributo che stabilisce se il messaggio contiene o meno una traccia video. Si ricorda che ogni messaggio contiene (di \inglese{default}) una traccia audio. Se tale attributo è impostato a \texttt{true} allora il messaggio contiene una traccia video.
	\item{\memberdata{-- date: Date}}
	Attributo che definisce l'orario di invio del messaggio.
\end{itemize}

\subsubsection*{Metodi}
\begin{description}
	\item{\method{+ getId(): Long}}\\
	Restituisce l'identificativo univoco del messaggio, ritornando il contenuto di \memberdata{id}.
	\item{\method{+ getSender(): IUserData}}\\
	Restituisce un istanza di tipo classname{abook.IUserData} che rappresenta il mittente del messaggio. Nello specifico il metodo ritorna l'attributo \memberdata{sender}.
	\item{\method{+ setSender(sender: IUserData): void}}\\
	Imposta il mittente del messaggio, sovrascrivendo il contenuto dell'attributo \memberdata{sender}.
	\item{\method{+ getReceiver(): IUserData}}\\
	Restituisce un istanza di tipo classname{abook.IUserData} che rappresenta il destinatario del messaggio. Nello specifico il metodo ritorna l'attributo \memberdata{receiver}.
	\item{\method{+ setReceiver(receiver: IUserData): void}}\\
	Imposta il destinatario del messaggio, sovrascrivendo il contenuto dell'attributo \memberdata{receiver}.
	\item{\method{+ isNew(): boolean}}\\
	Metodo che ritorna il contenuto dell'attributo \memberdata{status}.
	\item{\method{+ setNew(status: boolean): void}}\\
	Imposta il messaggio come ``già ascoltato'' o come ``da ascoltare'', sovrascrivendo il contenuto dell'attributo \memberdata{status}.
	\item{\method{+ isVideo(): boolean}}\\
	Restituisce il contenuto dell'attributo \memberdata{video}.
	\item{\method{+ setVideo(video: boolean): void }}\\	
	Metodo usato per impostare la ``natura'' del messaggio, impostando il contenuto dell'attributo \memberdata{video} mettendolo a \texttt{true} se il messaggio contiene una traccia video, oppure a \texttt{false} se non la contiene.
	\item{\method{+ getDate(): Date}}\\
	Restituisce la data in cui il mittente ha lasciato il messaggio nella segreteria del destinatario, ritornando il contenuto dell'attributo \memberdata{date}.
	\item{\method{+ setDate(date: Date): void}}\\
	Imposta la data in cui il mittente ha lasciato il messaggio nella segreteria del destinatario, sovrascrivendo il contenuto dell'attributo \memberdata{date}.

\end{description}

\subsection{Package org.softwaresynthesis.mytalk.server.authentication}\label{sec:authentication}

\classsection{ISecurityStrategy}

\subsubsection*{Funzione}
Interfaccia che identifica il comportamento di una strategia generica di crittografia dei dati.

\subsubsection*{Relazioni d'uso}
Nessuna relazione d'uso evidenziata.

\subsubsection*{Metodi}
\begin{description}
	\item{\method{+ encrypt(plainText: String): String}}\\
	Cripta la stringa di testo ricevuta come parametro.\\\\
	Il metodo può lanciare eccezioni:
	\begin{itemize}
		\item \exception{Exception}: il metodo lancia un'eccezione.
	\end{itemize}
	\item{\method{+ decrypt(encryptedText: String)}}\\
	Decripta la stringa di testo ricevuta come parametro.\\\\
	Il metodo può lanciare eccezioni:
	\begin{itemize}
		\item \exception{Exception}: il metodo lancia un'eccezione.
	\end{itemize}
\end{description}

\classsection{AESAlgorithm}

\subsubsection*{Funzione}
Implementazione della strategia di codifica/decodifica con l'uso dell'algoritmo AES a 128bit.

\subsubsection*{Relazioni d'uso}
\begin{itemize}
	\item \classname{ISecurityStrategy}: interfaccia d'implementazione
	\item \texttt{java.security.Key}: usata per creare un istanza di una chiave durante il processo di criptazione.
	\item \texttt{javax.crypto.Cipher}: usata per creare un istanza di un oggetto di criptazione che implementa l'algoritmo AES a 128 bit.
	\item \texttt{javax.crypto.spec.SecretKeySpec}: usata dall'algoritmo di per costruire una chiave segreta a partire da un array di byte.
	\item \texttt{sun.misc.BASE64Encoder}: utilizzata per eseguire una trasformazione da stringa a byte.
	\item \texttt{sun.misc.BASE64Decoder}: utilizzata per eseguire una trasformazione da byte in stringa.
\end{itemize}

\subsubsection*{Attributi}
\begin{description}
  \item{\memberdata{-- \{frozen\} \underline{keyValue}: byte[]}}\\
  Array di byte usato per definire il valore della chiave su cui si basa l'algoritmo AES di criptazione. La chiave effettiva sarà creata a partire da tale attributo, per mezzo della classe \texttt{javax.crypto.spec.SecretKeySpec}.
  \item{\memberdata{-- \{frozen\} \underline{algorithm}: String}}\\
  Stringa costante che identifica il nominativo dell'algoritmo usato, e che dovrà essere specificato durante l'uso di \texttt{javax.crypto.Cipher}. l'attributo avrà valore ``AES''.
\end{description}

\subsubsection*{Metodi}
\begin{description}
	\item{\method{- generateKey(): Key}}\\
	Metodo che restituisce una chiave di tipo \texttt{java.security.Key}, a partire dall'array \memberdata{keyValue}. Per fare ciò, il metodo usa il costruttore di \texttt{javax.crypto.spec.SecretKeySpec}.
	Il metodo può lanciare eccezioni:
	\begin{itemize}
		\item \exception{Exception}: il metodo può lanciare un'eccezione generica.
	\end{itemize}
	\item{\method{+ encrypt(plainText: String): String}}\\
	Metodo usato per criptare un testo di tipo \texttt{String} passato come parametro d'ingresso. Il metodo usa una chiave ottenuta a partire dal metodo \method{generateKey()} in associazione alla classe \texttt{javax.crypto.Cipher} per criptare il testo tramite l'algoritmo AES.
	\begin{itemize}
		\item \exception{Exception}: il metodo lancia un'eccezione.
	\end{itemize}
	\item{\method{+ decrypt(encryptedText: String)}}\\
	Decripta la stringa di testo ricevuta come parametro, attuando una procedura inversa a quella presentata nel metodo \method{encrypt(plainText: String)}\\\\
	Il metodo può lanciare eccezioni:
	\begin{itemize}
		\item \exception{Exception}: il metodo può lanciare un'eccezione generica.
	\end{itemize}
\end{description}

\classsection{PrincipalImpl}

\subsubsection*{Funzione}
Oggetto che permette di identificate univocamente uno \texttt{IUserData} memorizzato nel database del sistema \caName.

\subsubsection*{Relazioni d'uso}
\begin{itemize}
	\item \texttt{java.io.Serializable}: interfaccia d'implementazione usata per rendere serializzabili le istanze della classe.
	\item \texttt{java.security.Principal}: interfaccia d'implementazione usata per rendere caratterizzate le istanze della classe.
\end{itemize}

\subsubsection*{Attributi}
\begin{description}
  \item{\memberdata{-- \{frozen\} \underline{serialVersionUID}: long}}\\
  Identificativo univoco per la classe, usato al fine di rendere l'oggetto serializzabile.
  \item{\memberdata{-- mail: String}}\\
  Attributo che rappresenta l'indirizzo e-mail dell'utente.
\end{description}

\subsubsection*{Metodi}
\begin{description}
	\item{\method{+ PrincipalImpl(mail: String)}}\\
	Classe costruttore. crea un oggetto \texttt{PrincipalImpl} che permette di determinare univocamente lo \texttt{IUserData} che ha effettuato il login.
	\item{\method{+ getName(): String}}\\
	Restituisce l'elemento identificativo (mail dell'utente) dello \texttt{IUserData} che ha effettuato la procedura di login.
	\item{\method{+ equals(obj: Object): boolean}}\\
	Verifica l'uguaglianza di due oggetti \texttt{PrincipalImpl} sulla base di un confronto tra gli indirizzi mail degli utenti confrontati.
	\item{\method{+ hashCode(): int}}\\
	Restituisce il codice hash dell'oggetto di invocazione.				\item{\method{+ toString(): String}}\\
	Restituisce l'istanza dell'oggetto sotto forma di stringa. In particolare evidenziando l'indirizzo e-mail dell'utente.
\end{description}

\classsection{IAuthenticationData}

\subsubsection*{Funzione}
Interfaccia che identifica il comportamento di un oggetto adatto alla definizione dei dati di autenticazione di un utente.

\subsubsection*{Relazioni d'uso}
Nessuna relazione d'uso evidenziata.

\subsubsection*{Metodi}
\begin{description}
	\item{\method{+ getUsername(): String}}\\
Restituisce lo username fornito in input dall'utente durante la procedura di login.
	\item{\method{+ getPassword(): String}}\\
Restituisce la password fornita in input dall'utente durante la procedura di login
\end{description}

\classsection{AuthenticationData}

\subsubsection*{Funzione}
Oggetto contenente i dati di accesso utilizzati da un utente che vuole accedere al sistema \caName.

\subsubsection*{Relazioni d'uso}
\begin{itemize}
	\item \classname{IAuthenticationData}: interfaccia da implementare.
\end{itemize}

\subsubsection*{Attributi}
\begin{description}
  \item{\memberdata{-- username: String}}\\
	Attributo contenente l'username dell'utente che richiede l'autenticazione.
  \item{\memberdata{-- password: String}}\\
	Attributo contenente la password dell'utente che richiede l'autenticazione.
\end{description}

\subsubsection*{Metodi}
\begin{description}
	\item{\method{+ AuthenticationData(username: String, password: String)}}\\
	Costruttore pubblico che per creare un istanza di \texttt{AuthenticationData}, si basa sulla password e lo username dell'utente da autenticare.

	\item{\method{+ getUsername(): String}}\\
	Metodo che restituisce lo username dell'oggetto di autenticazione.
	
	\item{\method{+ getPassword(): String}}\\
	Metodo che Restituisce la password dell'utente che sta effettuando la procedura di login.

	\item{\method{+ hashCode(): int}}\\
	Restituisce il codice hash di questa istanza.
	
	\item{\method{+ equals(Object obj): boolean}}\\
	Determina l'uguaglianza di due istanze, effettuando un confronto tra l'oggetto stesso da cui è chiamato il metodo, e un Object \textit{obj} ricevuto come parametro d'ingresso.
	
	\item{\method{+ toString: String}}\\
	Metodo che Restituisce l'istanza sotto forma di stringa. Si osservi che il metodo non dovrà ritornare la password, ma solo lo username dell'utente da autenticare.

\end{description}

\classsection{AuthenticationModule}

%TODO definire questa parte tramite intervista a Mene.

\subsubsection*{Funzione}
Modulo di autenticazione utilizzato dal sistema \caName.

\subsubsection*{Relazioni d'uso}
\begin{itemize}
	\item \texttt{java.io.IOException}:
	\item \texttt{java.security.Principal}:
	\item \texttt{java.util.Map}:
	\item \texttt{java.util.Set}:
	\item \texttt{javax.security.auth.callback.Callback}:
	\item \texttt{javax.security.auth.callback.CallbackHandler}:
	\item \texttt{javax.security.auth.callback.NameCallback}:
	\item \texttt{javax.security.auth.callback.PasswordCallback}:
	\item \texttt{javax.security.auth.callback.UnsupportedCallbackException}:
	\item \texttt{javax.security.auth.login.FailedLoginException}:
	\item \texttt{javax.security.auth.login.LoginException}:
	\item \texttt{javax.security.auth.spi.LoginModule}:
	\item \texttt{javax.security.auth.Subject}:
	\item \classname{abook.IUserData}:
	\item \classname{dao.UserDataDAO}:
\end{itemize}

\subsubsection*{Attributi}
\begin{description}
  \item{\memberdata{-- login: boolean}}\\
  \item{\memberdata{-- commit: boolean}}\\
  \item{\memberdata{-- handler: CallbackHandler}}\\
  \item{\memberdata{-- password: char[]}}\\
  \item{\memberdata{-- username: String}}\\
  \item{\memberdata{-- principal: Principal}}\\
  \item{\memberdata{-- subject: Subject}}\\
\end{description}

\subsubsection*{Metodi}
\begin{description}
	\item{\method{+ initialize(subject: Subject, handler: CallbackHandler, sharedState: Map, option: Map): void }}\\

	\item{\method{+ login(): boolean throws LoginException}}\\
	
	\item{\method{+ commit(): boolean throws LoginException}}\\
	
	\item{\method{+ abort(): boolean throws LoginException}}\\
	
	\item{\method{+ logout(): boolean throws LoginException}}\\

\end{description}

\classsection{CredentialLoader}

\subsubsection*{Funzione}
Permette di caricare le credenziali di autenticazione, fornite dall'utente, per preparare la fase di login.

\subsubsection*{Relazioni d'uso}
\begin{itemize}
	\item \texttt{javax.security.auth.callback.Callback}: usato durante la procedura di inserimento dati.
	\item \texttt{javax.security.auth.callback.CallbackHandler}: Interfaccia implementata dalla classe.
	\item \texttt{javax.security.auth.callback.NameCallback}: tipo di dato richiesto in ingresso per completare la parte di login.
	\item \texttt{javax.security.auth.callback.PasswordCallback}: tipo di dato richiesto in ingresso per completare la parte di login.
	\item \texttt{java.io.IOException}: eccezione che può essere lanciata dal metodo \texttt{handle} definito dall'interfaccia \texttt{javax.security.auth.callback.CallbackHandler}.
	\item \texttt{javax.security.auth.callback.UnsupportedCallbackException}: eccezione che può essere lanciata dal metodo \texttt{handle} definito dall'interfaccia \texttt{javax.security.auth.callback.CallbackHandler}.
\end{itemize}

\subsubsection*{Attributi}
\begin{description}
  \item{\memberdata{-- credential: AuthenticationData}}\\
  Attributo contenente i dati da autenticare per la login dell'utente.
  \item{\memberdata{-- security: ISecurityStrategy}}\\
  Attributo che contiene un oggetto che definisce un algoritmo di criptazione per i dati. Necessario per criptare i dati di autenticazione.
\end{description}

\subsubsection*{Metodi}
\begin{description}
	\item{\method{+ CredentialLoader(credential: AuthenticationData, security: ISecurityStrategy)}}\\
	Costruttore pubblico. Crea un istanza con le credenziali fornite dall'utente (fornite in fase di login).

	\item{\method{+ handle(callbacks: Callback[]): void}}\\
	Effettua il caricamento e crittografa delle credenziali fornite dall'utente per la fase di login\\\\
	Il metodo può lanciare eccezioni:
	\begin{itemize}
		\item \exception{IOException}
		\item \exception{UnsupportedCallbackException}
	\end{itemize}

\end{description}

%TODO da rivedere nella definizione dei metodi! 
\classsection{LogoutServlet}

\subsubsection*{Funzione}
\inglese{Servlet} da richiamare per effettuare il logout dal sistema.

\subsubsection*{Relazioni d'uso}
\begin{itemize}
	\item \texttt{java.io.IOException}: eccezione richiamabile dai metodi \method{doPost()} e \method{doGet()}.
	\item \texttt{java.io.PrintWriter}: classe istanziata all'interno del metodo \method{doPost()}. Usata per scrivere l'output della \inglese{servlet}.
	\item \texttt{javax.servlet.ServletException}: eccezione rilancibile dai metodi \method{doPost()} e \method{doGet()}.
	\item \texttt{javax.servlet.http.HttpServlet}: classe estesa da \classname{LoginManager}.
	\item \texttt{javax.servlet.http.HttpServletRequest}:  classe per usata per la comunicazione con la \inglese{servlet} (inoltra i dati in ingresso alla \inglese{servlet}).
	\item \texttt{javax.security.auth.login.LoginContext}: classe usata in \method{doPost()} per eseguite la logout dal sistema.
	\item \texttt{javax.servlet.http.HttpServletResponse}: classe per usata per la comunicazione con la servlet (inoltra i dati in uscita dalla chiamata a servlet).
	\item \texttt{javax.servlet.http.HttpSession}: classe usata per definire una sessione HTTP.
	\item \classname{abook.IUserData}: usata per definire un utente.
	\item \classname{dao.UserDataDAO}: usata per comunicare tramite Hibernate con la tabella UserData della base di dati.
\end{itemize}

\subsubsection*{Attributi}
\begin{description}
  \item{\memberdata{\underline{-- {frozen} serialVersionUID: long}}}\\
  Utilizzato come ID univoco per identificare un oggetto serializzabile.
\end{description}

\subsubsection*{Metodi}
\begin{description}
	\item{\method{+ LogoutServlet()}}\\
	Costruttore della classe. Richiama il costruttore della classe padre \texttt{HttpServlet}.
	
	\item{\method{\# doGet(request: HttpServletRequest, response: HttpServletResponse): void}}\\
	Metodo di delega che richiama \method{doPost()} passando gli stessi dati in input.
	
	\item{\method{\# doPost(request: HttpServletRequest, response: HttpServletResponse): void}}\\
	Metodo che costituisce il kernel logico di risposta della \inglese{servlet}. Il metodo inizia caricando i parametri ricevuti mediante \texttt{HttpServletRequest}. Una chiamata a tale \inglese{servlet} corrisponde ad una richiesta di logout e si attua impostando a false la sessione contenuta nell'oggetto HttopServletRequest request. Il metodo procede con la creazione di un istanza  di \texttt{LoginContext} denominata context. Quindi viene effettuata la memorizzazione dell'oggetto ritornato da una chiamata:
\\
\verb|session.getAttribute("LoginContext");|
\\

L'oggetto cosi ottenuto dovrà essere controllato, ovvero si dovrà accertare che il tipo dinamico è conforme al tipo \texttt{LoginContext} (il programmatore dovrà obbligatoriamente usare   la primitiva \texttt{istanceof}). Quindi si sovrascrive il contenuto di context con quanto ottenuto dalla chiamata a metodo sopracitata, e si richiama il metodo \texttt{logout()} a partire dall'oggetto context. Il metodo termina invalidando la sessione richiamando il metodo \texttt{invalidate()} di HttpSession.

\end{description}

\classsection{LoginServlet}

\subsubsection*{Funzione}
\inglese{Servlet} da richiamare per effettuare l'autenticazione dell'utente.

\subsubsection*{Relazioni d'uso}
\begin{itemize}
	\item \texttt{java.io.IOException}: eccezione richiamabile dai metodi \method{doPost()} e \method{doGet()}.
	\item \texttt{java.io.PrintWriter}: classe istanziata all'interno del metodo \method{doPost()}. Usata per scrivere l'output della \inglese{servlet}.
	\item \texttt{javax.security.auth.login.LoginContext}: classe usata in \method{doPost()} per avviare la login dell'utente.
	\item \texttt{javax.servlet.ServletException}: eccezione richiamabile dai metodi \method{doPost()} e \method{doGet()}.
	\item \texttt{javax.servlet.http.HttpServlet}: classe estesa da \classname{LoginManager}.
	\item \texttt{javax.servlet.http.HttpServletRequest}:  classe usata per la comunicazione con la \inglese{servlet} (inoltra i dati in ingresso alla \inglese{servlet}).
	\item \texttt{javax.servlet.http.HttpServletResponse}: classe usata per la comunicazione con la \inglese{servlet} (inoltra i dati in uscita dalla chiamata a \inglese{servlet}).
	\item \texttt{javax.servlet.http.HttpSession}: classe usata per definire una sessione HTTP.
	\item \classname{abook.IUserData}: usata per definire un utente.
	\item \classname{dao.UserDataDAO}: usata per comunicare tramite Hibernate con la tabella \texttt{UserData} della base di dati.
\end{itemize}

\subsubsection*{Attributi}
\begin{description}
  \item{\memberdata{\underline{-- {frozen} serialVersionUID: long}}}\\
  Utilizzato come ID univoco per identificare un oggetto serializzabile.
\end{description}

\subsubsection*{Metodi}
\begin{description}
	\item{\method{+ LoginManager()}}\\
	Costruttore della classe. Richiama il costruttore della classe padre \texttt{HttpServlet}.
	
	\item{\method{\# doGet(request: HttpServletRequest, response: HttpServletResponse): void}}\\
	Metodo che richiama \method{doPost()} passando gli stessi dati in input.
	
	\item{\method{\# doPost(request: HttpServletRequest, response: HttpServletResponse): void}}\\

Metodo che costituisce il kernel logico di risposta della \inglese{servlet}. Il metodo inizia caricando i parametri ricevuti mediante \texttt{HttpServletRequest}. Una chiamata a tale metodo corrisponde ad una richiesta di login e si attua impostando due campi String con i dati ottenuti da una chiamata a \texttt{getParameter(``username'')} e \texttt{getParameter(``password'')} . Quindi si crea un istanza di \texttt{HttpSession} tramite il metodo \texttt{getSession(true)} del parametro request. Il metodo prosegue controllando se l'username e la password precedentemente ottenute hanno valore diverso da null. Nel caso il flusso principale continua creando un'istanza di \classname{AutenthicationData} (denominata credential) a partire dai parametri username e password. Quindi si carica la path del file di configurazione in un apposita stringa \texttt{pathFileConfig}, tramite la chiamata a metodo:\\
	
\verb|System.getenv("MyTalkConfiguration") + "\\LoginConfiguration.conf"|.\\

Il flusso principale prosegue all'interno di un blocco \texttt{try-catch} creando:
	\begin{itemize}
		\item[•] un \classname{CredentialLoader} (loader);
		\item[•] un \classname{LoginContext} (context);
		\item[•] un \classname{dao.UserDataDao} (user);
	\end{itemize}
Quindi tramite context si esegue il login e si impostano gli attributi di sessione come segue: 

\texttt{session.setAttribute("LoginContext", context);}

Infine si carica in user un istanza di \classname{dao.UserDataDAO} ottenuta dalla chiamata a metodo \method{UserDataDAO.getByEmail()}, e restituendo user in un formato di formattazione \texttt{Json} (\texttt{user.toJson()}).

\end{description}

\subsection{Package org.softwaresynthesis.mytalk.server.connection}\label{sec:connection}

\classsection{PushInbound}

\subsubsection*{Funzione}
Classe per la definizione di un apposito canale di comunicazione client-server. La classe è un'estensione di \texttt{org.apache.catalina.websocket.MessageInbound} che verrà usata poi in \classname{ChannelServlet}. Si osservi che l'associazione tra un istanza di \classname{PushInbound} e un utente del sistema è univoca: fintanto che connesso un utente ha un proprio \classname{PushInbound} presente sul server.

\subsubsection*{Relazioni d'uso}
\begin{itemize}
	\item \texttt{java.io.IOException}: Eccezione che può lanciare il metodo \method{OnTextMessage()}.
	\item \texttt{java.nio.CharBuffer}: tipologia di \textit{buffer} usata per memorizzare il messaggio ricevuto come parametro d'ingresso.
	\item \texttt{java.util.Iterator}: usata per scorrere il contenuto dell'insieme di \classname{AddressBookEntry}
	\item \texttt{java.util.Set}: struttura dati usata per memorizzare le \inglese{entry} (\classname{AddressBookEntry}) che costituiscono la rubrica dell'utente.
	\item \texttt{org.apache.catalina.websocket.MessageInbound}: classe da estendere, aggiunge alla classe \classname{PushInbound} le funzionalità necessarie per renderla un ``canale di comunicazione'' utilizzabile dai client.
	\item \texttt{com.google.gson.*}: converte i dati interni in una stringa formato json.
	\item \classname{State}: classe interna utilizzata per rappresentare lo stato dell'utente a cui è associata l'istanza di \classname{PushInbound}. Un utente ha uno stato fintanto che è connesso (quindi fintanto che esiste sul server un'istanza di \classname{PushInbound}  ad esso associata). Gli stati possibili sono: \texttt{available} (utente connesso e libero, quindi disposto a ricevere chiamate) e \texttt{occupied} (utente connesso ma attualmente occupato in un'altra conversazione).
	\item \classname{abook.AddressBookEntry}: usata per definire la rubrica dell'utente, richiamata dal metodo \method{onTextMessage} nel momento in cui si presenta la necessità di aggiornare lo stato dell'utente e renderlo visibile ai suoi contatti.
	\item \classname{abook.IUserData}: usata per riferire un istanza della classe ad un particolare utente.
	\item \classname{dao.UserDataDAO}: usata per riferire un istanza della classe ad un particolare utente.

\end{itemize}

\subsubsection*{Attributi}
\begin{description}
  \item{\memberdata{-- id: Long}}\\
  Identificativo di tipo \textit{Long} del canale di comunicazione associato ad un client univoco.
  \item{\memberdata{-- state: State}}\\
  Attributo usato per memorizzare lo stato dell'utente ``proprietario'' del \classname{PushInbound}.
\end{description}

%TODO da rivedere.

\subsubsection*{Metodi}
\begin{description}

	\item{\method{+ setId(n: Long): void}}\\
	Metodo per impostare il valore contenuto nell'attributo \memberdata{id}.
	
	\item{\method{+ getId(): Long}}\\
	Metodo che ritorna il contenuto dell'attributo \memberdata{id}.
	
	\item{\method{+ onTextMessage(message: CharBuffer): void}}\\
	Metodo invocato al momento della ricezione di un messaggio da parte del client. Il metodo riceve in input un oggetto di tipo \texttt{CharBuffer} contenente il messaggio inviato dal client.
Inizialmente il metodo crea un istanza per ognuno dei seguenti oggetti:
	\begin{itemize}
		\item[•]\texttt{Gson}: tipo di \textit{Json} definito da \textit{Google};
		\item[•]\texttt{JsonParser}: \textit{parser} usato per scorrere il contenuto di una stringa \textit{Json};
		\item[•]\texttt{JsonArray}: array popolato a partire dal \inglese{parsing} della stringa di messaggio data in input.
	\end{itemize}
	
	Prima di procedere si voglia considerare quanto segue: il metodo dopo aver ``segmentato'' il messaggio ricevuto in input, si occupa di esaminarne il contenuto che può essere di 5 tipologie, ciascuna identificata tramite un valore intero positivo da 1 a 5. Tale valore deve essere salvato nell'istanza di \texttt{JsonArray}
	Detto ciò, tornando a definire il flusso principale del metodo, si osservi che:

	\begin{itemize}
		\item[•]Se la richiesta inoltrata è del tipo 1: il metodo prende in lettura il messaggio e imposta il contenuto di \memberdata{id} con il valore letto mediante procedura di \inglese{parsing} analoga a quella definita al passo precedente. Per farlo utilizza il metodo \texttt{fromJson} richiamato dall'istanza \texttt{gsonObj} di tipo \texttt{Gson} creata al passo precedente. A tale metodo passa il contenuto dell'array e in particolare ciò che è salvato nella posizione 1 (utilizzo di metodo \texttt{get(int i)}).
		\item[•]Se la richiesta inoltrata è del tipo 2: il metodo procede con le istruzioni necessarie a scambiare i dati per la chiamata. Nello specifico viene salvato in un attributo di tipo \texttt{Long}, l'\texttt{id} del client che desidero chiamare, quindi ricerco l'oggetto \texttt{PushInbound} associato al client che desidero contattare, e gli inoltro il messaggio ricevuto come parametro d'ingresso.
		\item[•]Se la richiesta inoltrata è del tipo 3: il metodo comunica al client ``destinatario'' della chiamata, l'\texttt{id} del client chiamante. La procedura è analoga a quella identificata nel punto precedente, con la specifica che il messaggio inoltrato è l'identificativo del cliente che desidera avviare la chiamata.
		\item[•]Se la richiesta inoltrata è del tipo 4: il metodo si occupa dell'eliminazione della \textit{webSocket}, contattando l'istanza \classname{ChannelServlet} associata a questo \texttt{PushInbound}.
		\item[•]Se la richiesta inoltrata è del tipo 5: il metodo viene usato per modificare lo stato dell'utente, con il valore passato tramite messaggio. Dopo la ricezione del messaggio, il metodo cambia il valore del campo \texttt{state} con il valore ricevuto come parametro d'ingresso. Quindi procede ricavando la lista degli utenti nella rubrica dell'utente che ha cambiato stato e comunica loro che è avvenuto un cambiamento di stato.
	\end{itemize}
\end{description}

\classsection{ChanelServlet}

\subsubsection*{Funzione}
\inglese{Servlet} per la gestione delle connessioni. Richiamata dai client per ottenere le informazioni necessarie a stabilire una comunicazione client-client.

\subsubsection*{Relazioni d'uso}
\begin{itemize}
	\item \texttt{java.util.Vector}: vettore usato per memorizzare i \texttt{PushInBound} creati all'interno della classe.
	\item \texttt{javax.servlet.Servlet}: interfaccia d'implementazione. Necessaria per definire la classe come \inglese{servlet}.
	\item \texttt{javax.servlet.ServletConfig}: classe per la configura della \inglese{servlet}. 
	\item \texttt{javax.servlet.ServletException}: eccezione richiesta per la firma del metodo \method{init} (metodo richiesto per implementazione dell'interfaccia \texttt{javax.servlet.Servlet}.
	\item \texttt{javax.servlet.annotation.WebServlet}: necessaria per definire la \inglese{servlet} come \texttt{WebServlet}.
	\item \texttt{javax.servlet.http.HttpServletRequest}: necessaria per definire le richieste inoltrate alla classe (richieste a \inglese{servlet}).
	\item \classname{org.apache.catalina.websocket.StreamInbound}: classe utilizzata per definire un canale di comunicazione server-client.
	\item \classname{org.apache.catalina.websocket.WebSocketServlet}: classe da estendere per utilizzare i metodi di comunicazione tramite \inglese{webSocket}.
\end{itemize}

\subsubsection*{Attributi}
\begin{description}
  \item{\memberdata{\underline{-- \{frozen\} serialVersionUID: long}}}\\
  Utilizzato come ID univoco per identificare un oggetto serializzabile.
  \item{\memberdata{\underline{-- \{frozen\} clients: Vector<PushInbound>}}}\\
  Vettore contenente i canali di comunicazione server-client. Un canale di comunicazione \texttt{PushInbound} viene creato dal Server al momento dell'autenticazione di un utente (al quale poi viene associato).
\end{description}

\subsubsection*{Metodi}
\begin{description}
	\item{\method{+ ChannelServlet()}}\\
	Costruttore della \textit{servlet}. Richiama \method{super()}.
	
	\item{\method{\# createWebSocketInbound(subProtocol: String, request: HttpServletRequest): StreamInbound}}\\
	Metodo per la creazione di una \textit{websocket}. Il metodo ha anche il compito di salvare nel vettore di connessioni attive e la ritorna al client. Nello specifico il server segue l'iter:
	
	\begin{itemize}
		\item[1)] creazione di un \textit{PushInbound};
		\item[2)] aggiungo l'oggetto creato al punto precedente, al vettore di \textit{PushInbound} \memberdata{clients};
		\item[3)] il metodo termina ritornando l'istanza di \textit{PushInbound} creata al punto 1.
	\end{itemize}
	
	\item{\method{\underline{ + findClient(n: Long): PushInbound}}}\\
	Metodo per la ricerca di una connessione client dato l'identificativo dell'utente, esegue una ricerca (mediante ciclo \inglese{for}) all'interno del vettore \memberdata{clients} e ad ogni iterazione verifica se l'istanza i-esima presa in esame ha un campo id che corrisponde a quello dato in input come parametro di ricerca.
Il metodo termina ritornando il \textit{PushInbound} associato al client, se si verifica l'effettiva esitenza di un client con \texttt{id} uguale a quello passato in input, \texttt{null} altrimenti.
	
	\item{\method{\underline{ + removeClient(c: PushInbound): void}}}\\
	Metodo per la rimozione di un oggetto \texttt{PushInbound} dal vettore \memberdata{clients}. Il metodo esegue una ricerca tramite ciclo \inglese{for} (analoga a quella eseguita nel metodo \method{findClient}. Se la ricerca ha esito positivo il metodo termina eseguendo la rimozione dell'istanza dall'oggetto \memberdata{clients} tramite la chiamata a metodo \texttt{remove}.
\end{description}

\clearpage

\section{Specifica sotto-architettura clientpresenter}\label{sec:clientpresenterarchitecture}

La sotto-architettura \texttt{clientpresenter}, a livello di definizione di prodotto, richiede una trattazione speciale a causa del particolare dominio tecnologico coinvolto. Essa è infatti definita con il linguaggio \textit{JavaScript} che pur essendo definito come linguaggio orientato agli oggetti non permette di definire classi. A ciò si aggiunge il fatto che \textit{JavaScript} è debolmente tipizzato.
 In ragione di ciò, al fine di facilitare al programmatore la comprensione del progetto (pesando quindi alle varie entità come a classi) senza però confonderlo in fase di stesura del codice, si è pensato di adoperare la seguente terminologia.

\begin{description}
	\item{\scshape\bfseries Attributi}: saranno definiti con una sintassi simile a quella già usata. Quindi al nome dell'attributo sarà associato il tipo ``logico'' che idealmente rappresenta. In particolare emergono i seguenti casi:
		\begin{itemize}
			\item Nodi DOM: come si vedrà alcuni attributi rappresentano nodi DOM già esistenti nella pagina HTML che definisce l'interfaccia utente e modificati (a \textit{run-time}) sulla base di informazioni ottenute dal server. Tali attributi dovranno essere codificati come:
			
			\begin{verbatim}
				this.mio_attributo = document.getElementById(``id_DOM'');
			\end{verbatim}
			
			Tali attributi saranno indicati con il formalismo:
			
			\begin{verbatim}
				(+, -, #) nome_attributo: Nodo_DOM
			\end{verbatim}
			
			\item String: se l'attributo è destinato a contenere esclusivamente valori di tipo \texttt{String} sarà segnato nel documento come:
			
			\begin{verbatim}
				(+, -, #) nome_atributo: String
			\end{verbatim}
			
			\item Array: se l'attributo è un array, che in \textit{JavaScript} non chiede di essere definito per tipo di valori contenibili, nel documento sarà precisato come:
			
			\begin{verbatim}
				(+, -, #) nome_atributo: Tipo[]
			\end{verbatim}
		\end{itemize}
	\item{\scshape\bfseries Metodi}: la sintassi usata dal programmatore per definire un metodo in \textit{JavaScript} dovrà essere la seguente:
		\begin{verbatim}
			var nome_metodo = function(){ contenuto }
		\end{verbatim}
		
		mentre la sintassi usata nel documento per definirne la firma, resta la stessa usata in precedenza.
	
\end{description}

\subsection{Package org.softwaresynthesis.mytalk.clientpresenter.guicontrol}\label{sec:authentication}

%TODO da rivedere in modo BRUTALEEEEEEEEEEEEEEEEEEE!

\classsection{AddressBookPanelPresenter}

\subsubsection*{Funzione}
Presenter incaricato di gestire il pannello della rubrica, contiene le funzioni associate ai \inglese{widget} grafici della vista relativi a quest'ultima e ha la responsabilità di aggiornare la vista sulla base dei dati ricevuti dal server.

\subsubsection*{Relazioni d'uso}
Nessuna relazione d'uso evidenziata.

\subsubsection*{Attributi}
\begin{description}

  \item{\memberdata{-- element: DOM\_node}}\\
  Attributo che definisce il contenuto del nodo DOM inerente alla rappresentazione della rubrica. Nello specifico tale nodo corrisponde ad un \texttt{div} con id \texttt{AddressBookPanel}.
  \item{\memberdata{-- urServlet: String}}\\
  Attributo che definisce l'URL dal quale è richiamabile la \textit{servlet} usata per interagire con l'apparato server del sistema.
  \item{\memberdata{-- contacts: Array()}}\\
  Attributo che definisce la lista degli utenti presenti nella rubrica dell'utente che ha effettuato l'autenticazione.
   \item{\memberdata{-- groups: Array()}}\\
  Attributo che definisce la lista dei gruppi presenti nella rubrica dell'utente che ha effettuato l'autenticazione. Pero ogni gruppo sono  elencati i contatti contenuti in esso.
   \item{\memberdata{-- operations: String[]}}\\ 
   Array di stringhe usato per memorizzare i nomi delle istruzioni eseguibili. Tali nomi rappresentano una parte del nome della servlet incaricata a risolvere l'operazione descritta (e.g. GetContacts rappresenta l'operazione svolta dalla servlet avente nome: \memberdata{urlServlet} concatenato a ``GetContactsServlet''). I valori con i quali si dovrà caricare \memberdata{operations}) sono:
   \begin{itemize}
   	\item GetContacts;
    \item DoAddContact;
    \item DoDeleteContact;
    \item DoInsertInGroup;
    \item DoRemoveFromGroup;
    \item DoCreateGroup;
    \item DoDeleteGroup;
    \item GetGroups;
    \item DoBlock;
    \item DoUnblock;
    \item DoSearch;
   \end{itemize}

\end{description}

\subsubsection*{Metodi}
\begin{description}	
	
	\item{\method{+ inizialize(): void}}\\
	Metodo per inizializzare  \texttt{AddressBookPanel} e popolarlo con i contatti della rubrica. Il metodo modifica il DOM creando tre div: divSearch, divSort e divList. Quindi sequenzialmente ne crea il contenuto:
		\begin{itemize}
			\item divSearch: definisce un campo per l'input testuale e un bottone per ricercare nella lista un utente avente avete tra le parole chiave, una con valore uguale a quello inserito nel campo d'input.
			\item divSort: crea delle \inglese{select} per specificare le tipologie di ordinamento attuabili sulla lista degli utenti.
			\item divList: div che viene popolato con i nomi degli utenti presenti nella rubrica dell'utente che l'ha richiamata.
		\end{itemize}

Quindi dopo aver creato questi elementi il metodo procede inserendo nel DOM il ``sotto-albero'' cosi creato.

	\item{\method{+ getAddressBookContacts(): void}}\\
Metodo che recupera i contatti e i gruppi della propria rubrica dal server usando la tecnologia AJAX. Il metodo crea una variabile contenente un istanza \texttt{XMLHttpRequest()}. Quindi con metodo \inglese{POST} inoltra la richiesta alla \inglese{servlet} il cui indirizzo è generato dalla concatenazione di \memberdata{urlServlet} con \memberdata{operation[1]} seguito dal termine ``Servlet''. Quindi elabora i dati ottenuti mediante il metodo \texttt{JSON.parse()} e li salva negli array \memberdata{contacts} e \memberdata{groups}.
	
	\item{\method{+ setup(): void}}\\
	Inserisce i contatti estratti dal server all'interno della lista \texttt{AddressBookList} all'interno di \texttt{AddressBookPanel}.
	
	\item{\method{+ displayContactList(contact: String[]): void}}\\
	Metodo usato per visualizzare la lista dei contatti utente. Il metodo prende il nodo DOM (già esistente) e destinato a contenere la lista dei contatti, popolandolo richiamando il metodo \method{addListItem()}.
	
	\item{\method{-- addListItem(contact: String[]): void}}\\

Metodo richiamato da \method{dysplayContactList()} per visualizzare i contatti presenti all'array \memberdata{contacts} (passato come parametro). Il metodo crea le variabili contenenti i dati del contatto da attribuire al tag di markup \texttt{li}. Tali variabili sono quelle tipiche di un oggetti \classname{JSUserData} (nome, cognome, email, status e immagine).

	
	\item{\method{+ addContact(idContact: int): boolean}}\\
	Metodo per l'aggiunta di un contatto alla rubrica, riceve come parametro una stringa che definisce l'\texttt{id} dell'utente da inserire nella propria rubrica. 
Per assicurarsi di non introdurre ridondanza ed aggiungere un contatto già presente, il metodo esegue due controlli: uno lato client e uno lato server. 
Il controllo lato client, si scorre la lista \memberdata{contact} alla ricerca di un \texttt{id} che fa \inglese{match} con \texttt{idContact}. Nel caso già presente il metodo termina ritornando false. Per eseguire il controllo lato server il metodo contatta il server mediante la \inglese{servlet} \texttt{AddressBookDoAddContatctServlet}. Questa verificherà se l'utente avente \texttt{id} idContact è già presente nella rubrica, nel caso non lo sia  si occuperà di creare la nuova istanza di \texttt{AddressBookEntry} e la registrerà nel database, in caso contrario non esegue operazioni. 
La chiamata a servlet ritorna un valore booleano che sarà \texttt{true} se l'inserimento ha avuto successo, \texttt{false} altrimenti. 
Se l'inserimento ha avuto successo il flusso principale prosegue eseguendo l'aggiornamento della rubrica in locale tramite la chiamata a metodo \method{getAddressBookContacts()} seguita da \method{setup()}. Il metodo termina infine ritornando \texttt{true} se è stato aggiunto l'utente, \texttt{false} altrimenti.

	\item{\method{+ deleteContact(idContact: int): boolean}}\\
Metodo il cui funzionamento è analogo a quello precedentemente definito, si occupa di eliminare un contatto conoscendone l'\texttt{id} passato per parametro d'ingresso. Il flusso principale esegue due controlli per verificare la presenza dell'utente nella rubrica (analogamente a quanto visto nel caso precedente). 
Dopo tale verifica, se il contatto è realmente presente nella rubrica la \inglese{servlet} \texttt{AddressBookDoRemoveContactServlet} avrà eseguito la cancellazione come conseguenza al match effettuato in fase di verifica, nel caso il metodo richiama \method{getAddressBookContacts()} seguito da \method{setup()} per rieseguire l'aggiornamento della rubrica. Il metodo ritorna infine \texttt{true} se effettivamente il contatto era presente ed è stato eliminato, \texttt{false} altrimenti.

	\item{\method{+ addGroup(name: String): boolean}}\\
	Metodo per l'aggiunta di un nuovo gruppo nella rubrica, permette l'aggiunta di gruppi con nomi duplicati (ossia nomi di gruppi già presenti nella rubrica). Il flusso principale richiama la \inglese{servlet} \texttt{AddressBookDoInsertInGroupServlet} con il nome del gruppo da creare, quindi richiama prima \method{getAddressBookContacts()} seguito da \method{setup()} per aggiornare la rubrica in locale. Il metodo ritorna infine un \inglese{feedback} booleano per dare notifica sull'esito dell'operazione: \texttt{true} se è stato creato il nuovo gruppo, \texttt{false} altrimenti.
	
	\item{\method{+ deleteGroup(idGroup: int): boolean}}\\
	Metodo per l'eliminazione di un nuovo gruppo presente nella rubrica, riceve come parametro d'ingresso l'\texttt{id} del gruppo da eliminare. Il flusso principale richiama la \inglese{servlet} \texttt{AddressBookDoRemoveFromGroup} passandogli l'\texttt{id} del gruppo candidato all'eliminazione. Quindi, se la \inglese{servlet} notifica l'effettiva esistenza di un gruppo con tale \texttt{id}, essa si occupa dell'eliminazione e il flusso principale ritorna al metodo \method{deleteGroup} che richiama \method{getAddressBookContacts()} seguito da \method{setup()} per aggiornare la rubrica in locale. Il metodo ritorna un \inglese{feedback} booleano per dare notifica sull'esito dell'operazione: \texttt{true} se è stato creato il nuovo gruppo, \texttt{false} altrimenti.
	
	\item{\method{+ addContactInGroup(idContact: int, name: String): boolean}}\\
	Metodo per l'aggiunta di un contatto in un gruppo ben definito, riceve come parametro d'ingresso il nome del gruppo e l'\texttt{id} del contatto da aggiungervi. 
Il metodo richiama da prima \method{contactExistInGroup}, successivamente il flusso principale richiama la \inglese{servlet} \texttt{AddressBookDoInsertInGroupServlet} passandogli i parametri d'ingresso, quindi se la \inglese{servlet} notifica l'effettiva esistenza di un gruppo con il nome ricevuto, essa si occupa di aggiungervi il contatto e il flusso principale ritorna al metodo \method{addContactInGroup} che richiama \method{getAddressBookContacts()} seguito da  \method{setup()} per aggiornare la rubrica in locale. Il metodo ritorna infine un \inglese{feedback} booleano per dare notifica sull'esito dell'operazione: \texttt{true} se è stato creato il nuovo gruppo, \texttt{false} altrimenti.	
	
	\item{\method{+ deleteContactFromGroup(idContact: int, name: String): boolean}}\\
	Metodo per l'eliminazione di un contatto presente in un gruppo, riceve come parametro d'ingresso il nome del gruppo su cui effettuare la ricerca del contatto da eliminare e l'\texttt{id} di tale contatto. 
Il metodo richiama da prima \method{contactExistInGroup}, poi Il flusso principale richiama la \inglese{servlet} \texttt{AddressBookDoRemoveFromGroupServlet} passandogli i parametri d'ingresso, la \inglese{servlet} notifica l'effettiva esistenza di un gruppo con il nome ricevuto che contiene il contatto candidato all'eliminazione, essa si occupa dell'eliminazione e il flusso principale ritorna al metodo \method{deleteContactFromGroup} che richiama \method{getAddressBookContacts()} seguito da  \method{setup()} per aggiornare la rubrica in locale. Il metodo ritorna infine un \inglese{feedback} booleano per dare notifica sull'esito dell'operazione: \texttt{true} se è stato creato il nuovo gruppo, \texttt{false} altrimenti.
	
	\item{\method{+ contactExistInGroup(idContact: int, idGroup: int): boolean}}\\
Metodo che verifica l'esistenza di un contato in un determinato gruppo. Il metodo restituisce un feedback booleano per confermare o meno l'esistenza dell'utente ricercato nel gruppo designato.

\item{\method{+ applyFilterByString(param: String): array()}}\\
Metodo che data una stringa come parametro d'ingresso filtra la lista di contatti presente in locale. Inoltre non deve interagire con il server, ovvero non vengono richiamate \inglese{servlet}. 

In un primo step il programmatore dovrà creare un array() atto a contenere l'elenco dei contatti da restituire. Quindi il metodo procede con la definizione di un espressione regolare basata sul parametro ricevuto. Per tale operazione la sintassi obbligatoria è:\\

\verb|var pattern = new RegExp(param)|\\

Quindi viene avviato un ciclo for che analizza ogni contatto presente in \memberdata{contacts}. Successivamente si effettua un filtraggio sui campi name, surname, email del contatto attualmente in esame. Per fare tali operazioni si dovrà usare:\\

\verb|pattern.test(\*elemento su cui ricercare param*\)|\\

Se tale filtraggio ha esito positivo, il contatto viene aggiunto all'array dei contatti da restituire. Il metodo termina quando, dopo aver eseguito il ciclo for, il metodo ritorna l'array sopra-citato.

\item{\method{+ applyFilterByGroup(idGroup: int): array()}}\\
Metodo che dato un intero come parametro d'ingresso filtra la lista di contatti presente in locale, restituendo un elenco dei contatti appartenenti al gruppo avente id uguale a idGroup. Inoltre non deve interagire con il server, ovvero non vengono richiamate \inglese{servlet}. 

Come per \method{applyFilterByString()} anche in questo caso il metodo deve creare un Array() da popolare con i contatti appartenente al gruppo idGroup. Quindi viene eseguito un ciclo for che ispezione tutti i contatti presenti in \memberdata{contacts}. Per ogni ciclo il metodo controlla i gruppi in cui compare il contatto i-esimo (tale azione si attua per mezzo di un secondo ciclo innestato), e se si verifica che il gruppo in esame ha lo stesso id dell'identificativo idGroup, allora si aggiunge il contatto all'array creato in precedenza. Il metodo termina restituendo l'array sopra-citato.

\item{\method{+ showFilterp(filteredContacts: Array()): void}}\\
Metodo usato per l'eliminare il contenuto preesistente nella rubrica e visualizza il nuovo contenuto filtrato. Il metodo riceve come parametro un array che rappresenta la lista filtrata dei contatti da visualizzare. Per eseguire tali operazioni viene avviato un ciclo for per ogni contatto presente in \memberdata{contacts}. All'interno di tale ciclo deve essere eseguita la seguente istruzione:\\

\verb|addListItem(/*nome della variabile contenente il nodo dom AddressBookList*/,|\\
\verb|filtredContacts[contact])|\\

\end{description}

\classsection{LoginPanelPresenter}

\subsubsection*{Funzione}
Presenter incaricato di gestire il pannello di \textit{login}. prende in carica i dati inseriti dall'utente (username e password) passandoli al server per eseguire l'autenticazione. La \inglese{servlet} usata a tale scopo è \texttt{LoginManager}.

Il presenter si occupa anche di fornire delle funzionalità per il recupero della password.

\subsubsection*{Relazioni d'uso}

\subsubsection*{Attributi}
\begin{description}

  \item{\memberdata{-- element: DOM\_node}}\\
  Attributo che definisce il contenuto del nodo DOM inerente alla schermata di \inglese{login}, nello specifico tale nodo corrisponde ad un \texttt{<div>} con id \texttt{LoginPanel}.
  \item{\memberdata{-- urServlet: String}}\\
  Attributo contenente l'url della \inglese{servlet} incaricata di gestire il \inglese{login}.

\end{description}

\subsubsection*{Metodi}
\begin{description}

	\item{\method{-- testCredentials(data: String): void}}\\
Metodo che testa quanto ricevuto dal server e, in caso di login avvenuto correttamente, reindirizza il browser nel presenter finale dopo aver salvato i dati dell'utente. Il metodo inizia salvando in una variabile avente nome \inglese{user} il contenuto di data trattato attraverso una chiamata \texttt{JSON.parse(data)}, quindi se il contenuto di user non è nullo si procede richiamando il \texttt{communicationcenter} e impostando il valore di \texttt{communicationcenter.my} a \inglese{user}. Il metodo termina richiamando il metodo \method{hide()} e avviando la costruzione della UI mediante una chiamata \method{medietor.buildUI()}.
	
	\item{\method{+ sendAnswer(username: String, answer: String): void}}\\
	Metodo per inviare la risposta alla domanda segreta al server, riceve 2 parametri d'ingresso: lo username dell'utente e la risposta (answer) da inviare al server. Come primo passo viene creata un istanza di \texttt{XMLHttpRequest} salvata in una variabile avente nome request, successivamente si associa all'evento \texttt{onclick} di \textit{request} una funzione avente il compito di inviare la richiesta \textit{AJAX} al server, e di richiamare in risposta uno dei seguenti metodi:
	\begin{itemize}
		\item \method{correctAnswer()} nel caso la risposta inviata sia corretta;
		\item \method{incorrectAnswer()} se invece la risposta inviata dall'utente è errata.
	\end{itemize}
	
	\item{\method{+ correctAnswer(): void}}\\
	Metodo usato per modificare il nodo DOM salvato in \memberdata{element} per comunicare all'utente che la risposta da lui inserita (per il recupero password) è corretta. 
Il metodo deve rimuovere il figlio ``passwordretrieval'' da \memberdata{element} quindi salva in una variabile avente nome message il contenuto di una chiamata \verb|document.createElement("p")|. 
Successivamente procede caricando su message (usando il metodo \texttt{appendChild(var\_figlio)}) il messaggio da visualizzare per notificare all'utente il successo dell'operazione. Il messaggio visualizzato sarà: \\
	
	\verb|Recupero password avvenuto correttamente.|\\
	\verb|Ti è stata inviata un'email contenente i dati richiesti.|\\
	
il messaggio sopracitato è temporaneo e dovrà essere rimosso allo scadere di un timeout di durata 2000 millisecondi (usare la funzione \texttt{window.setTimeout(nome\_funzione))}.
	
	\item{\method{+ incorrectAnswer(): void}}\\
	Metodo che In caso di inserimento della risposta non corretta alla domanda segreta visualizza un messaggio di avvertimento all'utente per 2 secondi quindi lascia il controllo al \inglese{form} di inserimento della risposta alla domanda segreta.
	
	Il metodo deve cercare un elemento (\texttt{document.createElement("p")}) per la visualizzazione del messaggio:\\
	
	\verb|Dati non corretti. Inserire nuovamente la risposta.|\\
	
	Quindi si visualizza il messaggio con una procedura identica a quella usata in \method{correctAnswer()}.
		
	\item{\method{+ getSecretQuestion(): void}}\\
	Metodo che recupera la domanda segreta con una richiesta asincrona al server. Esso imposta una variabile \texttt{userID} con il valore ritornato da una chiamata \method{getUsername()}, quindi avvia una richiesta \textit{AJAX} inizializzando un oggetto \texttt{XMLHttpRequest}. La variabile \texttt{request} contenente tale istanza dovrà associare all'evento \texttt{onreadystatechange} una funzione che si occupa di catturare la stringa ritornata dal server (uso obbligatorio di \texttt{question = this.responseText}). Il metodo termina infine restituendo il valore contenuto in question.
	
	\item{\method{+ buildRetrievePasswordForm(): nodo\_DOM}}\\
	Metodo usato per costruire la \inglese{form} per il recupero della password utente, inizia costruendo la citata \inglese{form} con il valore ottenuto mediante una chiamata \verb|document.createElement("fieldset")|, quindi a partire dalla variabile in cui ha memorizzato il dato ritornato con il metodo precedente, imposta l'attributo ``id'' con ``passwordretrieval''. Definisce quindi una \inglese{label} con la domanda (utilizzare \verb|document.createElement("label")| e imposta l'attributo \texttt{for} con \textit{inputanswer}). Il metodo procede impostando il campo d'immissione della domanda e successivamente si impostano i seguenti attributi:
	\begin{itemize}
		\item attributo id con ``inputanswer'';
		\item attributo name con ``inputanswer'';
		\item attributo placeholder con ``risposta'';
		\item attributo required con ``required'';
	\end{itemize}
	
	Successivamente si crea il bottone di \texttt{submit} impostando gli attributi:
	\begin{itemize}
		\item attributo type con ``submit'';
		\item attributo value con ``OK'';
	\end{itemize}
	
	Associare all'evento \texttt{click} del bottone una funzione che richiami \method{sendAnswer()} passando come attributi \memberdata{username} e il valore contenuto nell'elemento \texttt{inputAnswer}.
	
	In ultima si usa \texttt{appendChild} per agganciare al nodo DOM la \inglese{label} e \texttt{inputAnswer}. Il metodo termina ritornando il nodo DOM.
	
	\item{\method{+ getUsername(): String}}\\
	Metodo  per il recupero dello username dall'interfaccia grafica. Per tale operazione il metodo deve usare l'istruzione:\\
	
	\verb|var username = document.getElementById("username").value;|
	\\
	
	quindi verifica se il contenuto di username è vuoto o no: nel caso lo sia, termina, altrimenti procede nel verificare se l'email inserita rispetta l'espressione regolare:\\
	
	\verb|^[A-Z0-9._%+-]+@[A-Z0-9.-]+\.[A-Z]{2,4}$|
	\\
	
	il metodo termina ritornando username.	
	
	\item{\method{+ getPassword(data: array()): String}}\\
	Metodo che esegue il \inglese{login} inviando al server i dati di autenticazione. Il metodo riceve come parametro d'ingresso un array associativo contenente \textit{username} e \textit{password}.
	
	Il flusso principale inizia verificando il che il contenuto di \texttt{data.username} e di \texttt{data.password} non sia vuoto, quindi invia una richiesta \textit{AJAX} al server (come sempre in questo caso è d'obbligo usare \verb|var request = new XMLHttpRequest();|).
	Il programmatore dovrà associare all'evento \texttt{onreadystatechange} una funzione che si occupa di richiamare il metodo \method{testCredentials()}. Il passo successivo consiste ``nell'aprire'' la richiesta request con un istruzione open() che riceve come parametro l'url della \inglese{servlet} associata (valore presente in \memberdata{servletURL}), successivamente viene usata un istruzione \texttt{setRequestHeader()} passando i seguenti parametri:
	\begin{itemize}
		\item[•] ``content-type'';
		\item[•] ``application/x-www-form-urlencoded'';
	\end{itemize}
	infine va creata la \textit{querystring} da inviare (tramite istruzione request.send()). Tale \inglese{query} deve essere obbligatoriamente impostata a:\\
	
	\verb|"username=" + encodeURIComponent(data.username) +|\\ \verb|"&password=" + encodeURIComponent(data.password) + "&operation=1"|\\
	
	Il metodo termina ritornando la \textit{querystring}.

	\item{\method{+ hide(): void}}\\
	Metodo usato per nascondere il \inglese{form} di autenticazione, per lasciare spazio nella finestra ad altri elementi grafici come la schermata principale o il pannello di registrazione.
	
	\item{\method{+ initialize(): void}}\\
	Metodo usato per inizializzazione del pannello di login con la creazione di tutti i \inglese{widget} grafici che sono contenuti al suo interno. Nello specifico il metodo esegue la creazione (nell'ordine proposto) dei seguenti elementi:
	\begin{itemize}
		\item form, con attributi:
		\begin{itemize}
			\item ``name'' impostato a ``login'';
			\item ``action'' impostato a ``'';
			\item ``method'' impostato a ``'';
			\item ``accept-charset'' impostato a ``utf-8'';				
		\end{itemize}
		\item ul contennete il form;
		\item li che conterrà lo username, con attributi;
		\item label che riporterà la frase ``Nome utente: ''. Label con attributi:
		\begin{itemize}
			\item ``for'' impostato a ``username'';		
		\end{itemize}
		\item input per l'inserimento della e-mail, con attributi:
		\begin{itemize}
			\item ``type'' impostato a ``email'';
			\item ``id'' impostato a ``username'';
			\item ``name'' impostato a ``username'';
			\item ``placeholder'' impostato a ``yourname@email.com'';				\item ``required'' impostato a ``required'';
		\end{itemize}		
		\item input per l'inserimento della password, con attributi:
		\begin{itemize}
			\item ``type'' impostato a ``password'';
			\item ``id'' impostato a ``password'';
			\item ``name'' impostato a ``password'';
			\item ``placeholder'' impostato a ``password'';				\item ``required'' impostato a ``required'';
		\end{itemize}		
		\item input submit pulsante di login, con attributi:
		\begin{itemize}
			\item ``type'' impostato a ``submit'';
			\item ``value'' impostato a ``Login'';
		\end{itemize}
		\item input submit pulsante di registrazione, con attributi:
		\begin{itemize}
			\item ``type'' impostato a ``submit'';
			\item ``value'' impostato a ``Registrazione'';
		\end{itemize}
		\item input submit pulsante di recupero password, con attributi:
		\begin{itemize}
			\item ``type'' impostato a ``submit'';
			\item ``value'' impostato a ``Recupera password'';
		\end{itemize}
	\end{itemize}

\end{description}

\classsection{RegisterPanelPresenter}

\subsubsection*{Funzione}
Presenter incaricato di gestire la \inglese{form} di registrazione

\subsubsection*{Relazioni d'uso}

Nessuna relazione d'uso evidenziata.

\subsubsection*{Attributi}
\begin{description}

  \item{\memberdata{-- servletURL: String}}\\
  Url della \inglese{servlet} che deve gestire la registrazione.
  \item{\memberdata{-- element: DOM\_node}}\\
  Elemento (inteso come nodo DOM) controllato da questo Presenter.

\end{description}

\subsubsection*{Metodi}
\begin{description}

	\item{\method{+ getSurname(): String}}\\
	Metodo che estrae dalla \inglese{form} il valore del cognome del nuovo utente e lo restituisce.
	
	\item{\method{+ getName(): String}}\\
	Metodo che estrae dalla \inglese{form} il valore del nome del nuovo utente e lo restituisce.
	
	\item{\method{+ getAnswer(): String}}\\
	Metodo che estrae dal \inglese{form} la risposta alla domanda segreta associata al nuovo utente e la restituisce.	
		
	\item{\method{+ getQuestion(): String}}\\
	Metodo che estrae dalla \inglese{form} la domanda segreta associata al nuovo utente e la restituisce.
	
	\item{\method{+ getPassword(): String}}\\
	Metodo che estrae dalla \inglese{form} la password associata al nuovo utente e la restituisce.
	
	\item{\method{+ getUsername(): String}}\\
	Metodo che estrae dalla \inglese{form} lo username del nuovo utente (per username si intende l'indirizzo email dell'utente). Si osservi che il metodo dovrà verificare che la mail restituita sia ben formattata secondo l'espressione regolare:\\
	
	\verb|^[A-Z0-9._%+-]+@[A-Z0-9.-]+\.[A-Z]{2,4}$|
	\\
	
	Se la mail è valida il metodo termina ritornandone il valore.
	
	\item{\method{+ getPicturePath(): String}}\\
	Metodo che estrae dalla \inglese{form} il percorso dell'immagine del nuovo profilo utente e lo restituisce.
	
	\item{\method{+ register(userData: Arrey()): String}}\\
	Metodo che Invia i dati ricevuti alla \inglese{servlet} per la creazione di un nuovo account utente. Tale metodo inizia creando un istanza (\inglese{request}) di \texttt{XMLHttpRequest} usata per comunicare con la \inglese{servlet}. Quindi il metodo abilita \textit{request} ad inviare un richiesta sincrona al server, usando l'istruzione:\\
	
	\verb|request.open("POST", servletURL, false)|\\
	
	Il metodo procede definendo la stringa da usare come \textit{query} per da inviare alla \inglese{servlet}. Tale \textit{querystring} contiene:\\
	
	\verb|"username=" + encodeURIComponent(userData.username) +|\\
	\verb|"&password=" + encodeURIComponent(userData.password) +|\\
	\verb|"&question=" + encodeURIComponent(userData.question) +|\\
	\verb|"&answer=" + encodeURIComponent(userData.answer)|\\
	
	Il metodo termina restituendo il contenuto di \textit{querystring}.
	
	\item{\method{+ initialize(): void}}\\
	Metodo che si occupa di inizializzare la \inglese{form} di registrazione con la creazione di tutti i \inglese{widget} grafici che sono contenuti al suo interno. Nello specifico il metodo esegue i seguenti passi:
	\begin{itemize}
		\item[•] creazione dell'elemento \inglese{form};
		\item[•] creazione dell'elemento <ul> contenuto nel \inglese{form};		
		\item[•] creazione dell'item per lo username, definito mediante una \inglese{label} riportante la frase ``Indirizzo email: '', e di un campo input di tipo email per l'inserimento della email dell'utente;
		\item[•] costruisce il \inglese{list item} con la \inglese{label} e l'input definiti al punto precedente, usando il metodo \texttt{appendChild()};
		\item[•] crea l'item per la password, definendo una \inglese{label} riportante la frase ``Password: '' e un oggetto input di tipi password;			
		\item[•] costruisce il \inglese{list item} con la \inglese{label} e l'input definiti al punto precedente, usando il metodo \texttt{appendChild()};
		\item[•] creazione dell'item per la domanda segreta definito mediante una \inglese{label} avente valore ``Domanda segreta: '' e un oggetto input di tipo \texttt{text};
		\item[•] costruisce il \inglese{list item} con la \inglese{label} e l'input definiti al punto precedente, usando il metodo \texttt{appendChild()};
		\item[•] creazione dell'item per la risposta alla domanda segreta, definito mediante una \inglese{label} avente valore ``Risposta: '' e un oggetto input di tipo \texttt{text};
		\item[•] costruisce il \inglese{list item} con la \inglese{label} e l'input definiti al punto precedente, usando il metodo \texttt{appendChild()};
		\item[•] creazione dell'item per il nome, definito mediante una \inglese{label} avente valore ``Nome: '' e un oggetto input di tipo text;
		\item[•] costruisce il \inglese{list item} con la \inglese{label} e l'input definiti al punto precedente, usando il metodo \texttt{appendChild()};
		\item[•] creazione dell'item per il cognome, definito mediante una \inglese{label} avente valore ``Cognome: '' e un oggetto input di tipo \texttt{text};
		\item[•] costruisce il \inglese{list item} con la \inglese{label} e l'input definiti al punto precedente, usando il metodo \texttt{appendChild()};
		\item[•] creazione dell'item per inserire l'immagine utente, definito mediante una \inglese{label} avente valore ``Immagine del profilo: '' e un oggetto input di tipo \texttt{text};
		\item[•] costruisce il \inglese{list item} con la \inglese{label} e l'input definiti al punto precedente, usando il metodo \texttt{appendChild()};
		\item[•] creazione di un pulsante di registrazione (oggetto input di tipo submit e value ``Registrati'').
		\item[•] operazione finale per ``appendere'' tutti  gli oggetti creati al \texttt{registerForm}.
	\end{itemize}
	
	\item{\method{+ hide(): void}}\\
	Nasconde il \inglese{form} di registrazione per lasciare spazio alla schermata principale dell'applicativo (che deve essere costruita dal PresenterMediator).

\end{description}

\clearpage

\section{Specifica sotto-architettura clientview}\label{sec:clientviewarchitecture}

\clearpage

\section{Tracciamento della relazione componenti-requisiti}

\end{document}
